
\documentclass{article}
\usepackage{colortbl}
\usepackage{makecell}
\usepackage{multirow}
\usepackage{supertabular}

\begin{document}

\newcounter{utterance}

\centering \large Interaction Transcript for game `clean\_up', experiment `1\_medium\_5obj\_cs', episode 1 with gemma{-}4{-}26b{-}a4b{-}it.
\vspace{24pt}

{ \footnotesize  \setcounter{utterance}{1}
\setlength{\tabcolsep}{0pt}
\begin{supertabular}{c@{$\;$}|p{.15\linewidth}@{}p{.15\linewidth}p{.15\linewidth}p{.15\linewidth}p{.15\linewidth}p{.15\linewidth}}
    \# & $\;$A & \multicolumn{4}{c}{Game Master} & $\;\:$B\\
    \hline

    \theutterance \stepcounter{utterance}  
    & & \multicolumn{4}{p{0.6\linewidth}}{
        \cellcolor[rgb]{0.9,0.9,0.9}{
            \makecell[{{p{\linewidth}}}]{
                \texttt{\tiny{[P1$\langle$GM]}}
                \texttt{Jsem váš herní mistr a hrajete společnou hru s následující mřížkou jako herním plánem:} \\
\\ 
\texttt{```} \\
\texttt{ 1234567} \\
\texttt{╔══╤══╤═╗} \\
\texttt{║◌◌│T◌│◌║ 1} \\
\texttt{╟──┼──┴─╢ 2} \\
\texttt{║WI│◌◌◌◌║ 3} \\
\texttt{╟──┘◌◌◌◌║ 4} \\
\texttt{║◌◌◌H◌◌C║ 5} \\
\texttt{║◌◌┌─┬──╢ 6} \\
\texttt{║◌◌│◌│◌◌║ 7} \\
\texttt{╚══╧═╧══╝} \\
\\ 
\texttt{```} \\
\\ 
\texttt{* Horní okraj zobrazuje souřadnice x rostoucí doprava a pravý okraj souřadnice y rostoucí dolů.} \\
\texttt{* Na vaší mřížce jsou náhodně umístěny následující objekty: 'W', 'I', 'T', 'C', 'H'.} \\
\\ 
\\ 
\texttt{Druhý hráč vidí variantu herní desky, kde jsou objekty umístěny na různých náhodných místech. Vy nevidíte desku druhého hráče a on nevidí tu vaši.} \\
\\ 
\\ 
\texttt{**Cíl:**} \\
\\ 
\texttt{Oba hráči musí přesunout objekty na svém pozadí tak, aby stejné objekty skončily na stejných souřadnicích. Musíte komunikovat s druhým hráčem, abyste se dohodli na společné cílové konfiguraci.} \\
\\ 
\\ 
\texttt{**Pravidla:**} \\
\\ 
\texttt{* V každém tahu můžete poslat přesně jeden z následujících dvou příkazů:} \\
\texttt{ 1. `ŘEKNI: <ZPRÁVA>`: odeslání zprávy (vše až do dalšího zalomení řádku) druhému hráči. Přepošlu ji vašemu partnerovi.} \\
\texttt{2. `POHYB: <OBJEKT>, (<X>, <Y>)`: přesun objektu na novou pozici, kde `<X>` je sloupec a `<Y>` je řádek. Budu vás informovat, zda byl váš přesun úspěšný, nebo ne.} \\
\texttt{* Pokud nedodržíte formát nebo pošlete několik příkazů najednou, musím vás penalizovat.} \\
\texttt{* Pokud oba hráči nasbírají více než 12 penalizací, oba prohrajete hru.} \\
\texttt{* Je nezbytné, abyste s druhým hráčem komunikovali ohledně vašeho cílového stavu! *Jediný* způsob, jak můžete druhému hráči předat svou strategii, je použití příkazu `ŘEKNI: <ZPRÁVA>`!} \\
\\ 
\\ 
\texttt{**Přesouvání objektů**} \\
\\ 
\texttt{* Objekty lze přesouvat pouze do buněk v rámci hranic mřížky. Cílová buňka musí být prázdná, tj. musí obsahovat pouze symbol '◌'.} \\
\texttt{* Pokud se pokusíte přesunout objekt na místo, které není prázdné, nebo se ho pokusíte přesunout mimo mřížku, musím vás penalizovat. Dostanete další pokus.} \\
\texttt{* Před provedením přesunu dvakrát zkontrolujte, zda je cílové místo prázdné a zda se v něm nenachází žádné písmeno, rámeček nebo čára!} \\
\\ 
\\ 
\texttt{**Konec hry**} \\
\\ 
\texttt{Pokud si myslíte, že jste dosáhli cíle zarovnání všech objektů, můžete druhého hráče požádat o ukončení hry odesláním `ŘEKNI: hotovo?`. Pokud vás druhý hráč požádá o ukončení hry a vy odpovíte `ŘEKNI: hotovo!`, hra skončí.} \\
\\ 
\texttt{Oba hráči vyhrají, pokud hra skončí během 20 kol, přičemž jedno kolo je definováno jako odeslání platného příkazu oběma hráči.} \\
\\ 
\texttt{**Bodování:**} \\
\\ 
\texttt{Čím blíže jsou stejné objekty na obou herních deskách, tím více bodů získáte. Penalizace vám body snižují. Podaří se vám překonat rekord?} \\
\\ 
\\ 
\texttt{Pošlete prosím zprávu druhému hráči a začněte hru!} \\
            }
        }
    }
    & \\ \\

    \theutterance \stepcounter{utterance}  
    & \multicolumn{4}{p{0.6\linewidth}}{
        \cellcolor[rgb]{0.9,0.9,0.9}{
            \makecell[{{p{\linewidth}}}]{
                \texttt{\tiny{[P1$\rangle$GM]}}
                \texttt{ŘEKNI: Ahoj! Musíme najít společné souřadnice pro všechny objekty (W, I, T, C, H). Já vidím jejich aktuální pozice: W(1,3), I(2,3), T(4,1), H(4,5), C(7,5). Navrhuji, abychom se pokusili posunout všechny objekty do středové části nebo do nějaké jednoduché konfigurace. Jak vypadá tvá deska?} \\
            }
        }
    }
    & & \\ \\

    \theutterance \stepcounter{utterance}  
    & & & \multicolumn{2}{p{0.3\linewidth}}{
        \cellcolor[rgb]{0.9,0.9,0.9}{
            \makecell[{{p{\linewidth}}}]{
                \texttt{\tiny{[GM$|$GM]}}
                \texttt{ŘEKNI: Ahoj! Musíme najít společné souřadnice pro všechny objekty (W, I, T, C, H). Já vidím jejich aktuální pozice: W(1,3), I(2,3), T(4,1), H(4,5), C(7,5). Navrhuji, abychom se pokusili posunout všechny objekty do středové části nebo do nějaké jednoduché konfigurace. Jak vypadá tvá deska?} \\
            }
        }
    }
    & & \\ \\

    \theutterance \stepcounter{utterance}  
    & & \multicolumn{4}{p{0.6\linewidth}}{
        \cellcolor[rgb]{0.9,0.9,0.9}{
            \makecell[{{p{\linewidth}}}]{
                \texttt{\tiny{[GM$\rangle$P2]}}
                \texttt{Jsem váš herní mistr a hrajete společnou hru s následující mřížkou jako herním plánem:} \\
\\ 
\texttt{```} \\
\texttt{ 1234567} \\
\texttt{╔══╤══╤═╗} \\
\texttt{║◌◌│W◌│I║ 1} \\
\texttt{╟──┼──┴─╢ 2} \\
\texttt{║◌◌│◌◌◌◌║ 3} \\
\texttt{╟──┘T◌◌◌║ 4} \\
\texttt{║H◌◌◌◌◌◌║ 5} \\
\texttt{║◌◌┌─┬──╢ 6} \\
\texttt{║C◌│◌│◌◌║ 7} \\
\texttt{╚══╧═╧══╝} \\
\\ 
\texttt{```} \\
\\ 
\texttt{* Horní okraj zobrazuje souřadnice x rostoucí doprava a pravý okraj souřadnice y rostoucí dolů.} \\
\texttt{* Na vaší mřížce jsou náhodně umístěny následující objekty: 'W', 'I', 'T', 'C', 'H'.} \\
\\ 
\\ 
\texttt{Druhý hráč vidí variantu herní desky, kde jsou objekty umístěny na různých náhodných místech. Vy nevidíte desku druhého hráče a on nevidí tu vaši.} \\
\\ 
\\ 
\texttt{**Cíl:**} \\
\\ 
\texttt{Oba hráči musí přesunout objekty na svém pozadí tak, aby stejné objekty skončily na stejných souřadnicích. Musíte komunikovat s druhým hráčem, abyste se dohodli na společné cílové konfiguraci.} \\
\\ 
\\ 
\texttt{**Pravidla:**} \\
\\ 
\texttt{* V každém tahu můžete poslat přesně jeden z následujících dvou příkazů:} \\
\texttt{ 1. `ŘEKNI: <ZPRÁVA>`: odeslání zprávy (vše až do dalšího zalomení řádku) druhému hráči. Přepošlu ji vašemu partnerovi.} \\
\texttt{2. `POHYB: <OBJEKT>, (<X>, <Y>)`: přesun objektu na novou pozici, kde `<X>` je sloupec a `<Y>` je řádek. Budu vás informovat, zda byl váš přesun úspěšný, nebo ne.} \\
\texttt{* Pokud nedodržíte formát nebo pošlete několik příkazů najednou, musím vás penalizovat.} \\
\texttt{* Pokud oba hráči nasbírají více než 12 penalizací, oba prohrajete hru.} \\
\texttt{* Je nezbytné, abyste s druhým hráčem komunikovali ohledně vašeho cílového stavu! *Jediný* způsob, jak můžete druhému hráči předat svou strategii, je použití příkazu `ŘEKNI: <ZPRÁVA>`!} \\
\\ 
\\ 
\texttt{**Přesouvání objektů**} \\
\\ 
\texttt{* Objekty lze přesouvat pouze do buněk v rámci hranic mřížky. Cílová buňka musí být prázdná, tj. musí obsahovat pouze symbol '◌'.} \\
\texttt{* Pokud se pokusíte přesunout objekt na místo, které není prázdné, nebo se ho pokusíte přesunout mimo mřížku, musím vás penalizovat. Dostanete další pokus.} \\
\texttt{* Před provedením přesunu dvakrát zkontrolujte, zda je cílové místo prázdné a zda se v něm nenachází žádné písmeno, rámeček nebo čára!} \\
\\ 
\\ 
\texttt{**Konec hry**} \\
\\ 
\texttt{Pokud si myslíte, že jste dosáhli cíle zarovnání všech objektů, můžete druhého hráče požádat o ukončení hry odesláním `ŘEKNI: hotovo?`. Pokud vás druhý hráč požádá o ukončení hry a vy odpovíte `ŘEKNI: hotovo!`, hra skončí.} \\
\\ 
\texttt{Oba hráči vyhrají, pokud hra skončí během 20 kol, přičemž jedno kolo je definováno jako odeslání platného příkazu oběma hráči.} \\
\\ 
\texttt{**Bodování:**} \\
\\ 
\texttt{Čím blíže jsou stejné objekty na obou herních deskách, tím více bodů získáte. Penalizace vám body snižují. Podaří se vám překonat rekord?} \\
\\ 
\texttt{Druhý hráč zahájil hru odesláním této zprávy:} \\
\\ 
\texttt{"Ahoj! Musíme najít společné souřadnice pro všechny objekty (W, I, T, C, H). Já vidím jejich aktuální pozice: W(1,3), I(2,3), T(4,1), H(4,5), C(7,5). Navrhuji, abychom se pokusili posunout všechny objekty do středové části nebo do nějaké jednoduché konfigurace. Jak vypadá tvá deska?"} \\
\\ 
\texttt{Jaký je váš první příkaz?} \\
            }
        }
    }
    & \\ \\

    \theutterance \stepcounter{utterance}  
    & & & \multicolumn{4}{p{0.6\linewidth}}{
        \cellcolor[rgb]{0.9,0.9,0.9}{
            \makecell[{{p{\linewidth}}}]{
                \texttt{\tiny{[GM$\langle$P2]}}
                \texttt{ŘEKNI: Ahoj! Moje deska vypadá takto: W(4,1), I(7,1), T(4,4), C(1,7), H(1,5). Navrhuji, abychom se pokusili posunout všechny objekty do řady nebo sloupce, aby to bylo snadno koordinovatelné. Co říkáš na to, aby všechny objekty skončily na souřadnicích (1,1) až (5,1)? Tedy: W(1,1), I(2,1), T(3,1), C(4,1), H(5,1).} \\
            }
        }
    }
     \\ \\

    \theutterance \stepcounter{utterance}  
    & & & \multicolumn{2}{p{0.3\linewidth}}{
        \cellcolor[rgb]{0.9,0.9,0.9}{
            \makecell[{{p{\linewidth}}}]{
                \texttt{\tiny{[GM$|$GM]}}
                \texttt{ŘEKNI: Ahoj! Moje deska vypadá takto: W(4,1), I(7,1), T(4,4), C(1,7), H(1,5). Navrhuji, abychom se pokusili posunout všechny objekty do řady nebo sloupce, aby to bylo snadno koordinovatelné. Co říkáš na to, aby všechny objekty skončily na souřadnicích (1,1) až (5,1)? Tedy: W(1,1), I(2,1), T(3,1), C(4,1), H(5,1).} \\
            }
        }
    }
    & & \\ \\

    \theutterance \stepcounter{utterance}  
    & & \multicolumn{4}{p{0.6\linewidth}}{
        \cellcolor[rgb]{0.9,0.9,0.9}{
            \makecell[{{p{\linewidth}}}]{
                \texttt{\tiny{[P1$\langle$GM]}}
                \texttt{Vaše zpráva byla předána druhému hráči.} \\
\\ 
\\ 
\texttt{Právě hrajete kolo 1 z maximálně 20.} \\
\texttt{Dohromady jste nasbírali 0 z 12 pokut.} \\
\texttt{Druhý hráč poslal tuto zprávu:} \\
\\ 
\\ 
\texttt{"Ahoj! Moje deska vypadá takto: W(4,1), I(7,1), T(4,4), C(1,7), H(1,5). Navrhuji, abychom se pokusili posunout všechny objekty do řady nebo sloupce, aby to bylo snadno koordinovatelné. Co říkáš na to, aby všechny objekty skončily na souřadnicích (1,1) až (5,1)? Tedy: W(1,1), I(2,1), T(3,1), C(4,1), H(5,1)."} \\
\\ 
\\ 
\texttt{Jaký je váš další příkaz?} \\
            }
        }
    }
    & \\ \\

    \theutterance \stepcounter{utterance}  
    & \multicolumn{4}{p{0.6\linewidth}}{
        \cellcolor[rgb]{0.9,0.9,0.9}{
            \makecell[{{p{\linewidth}}}]{
                \texttt{\tiny{[P1$\rangle$GM]}}
                \texttt{ŘEKNI: To je skvělý plán! Souřadnice (1,1) až (5,1) jsou jasné. Budu se snažit posunout své objekty na tyto pozice. Začnu s objektem T, který je už blízko.} \\
\\ 
\texttt{POHYB: T, (3,1)} \\
            }
        }
    }
    & & \\ \\

    \theutterance \stepcounter{utterance}  
    & & & \multicolumn{2}{p{0.3\linewidth}}{
        \cellcolor[rgb]{0.9,0.9,0.9}{
            \makecell[{{p{\linewidth}}}]{
                \texttt{\tiny{[GM$|$GM]}}
                \texttt{ŘEKNI: To je skvělý plán! Souřadnice (1,1) až (5,1) jsou jasné. Budu se snažit posunout své objekty na tyto pozice. Začnu s objektem T, který je už blízko.} \\
\\ 
\texttt{POHYB: T, (3,1)} \\
            }
        }
    }
    & & \\ \\

    \theutterance \stepcounter{utterance}  
    & & & \multicolumn{2}{p{0.3\linewidth}}{
        \cellcolor[rgb]{0.9,0.9,0.9}{
            \makecell[{{p{\linewidth}}}]{
                \texttt{\tiny{[GM$|$GM]}}
                \texttt{Invalid response format: ŘEKNI: To je skvělý plán! Souřadnice (1,1) až (5,1) jsou jasné. Budu se snažit posunout své objekty na tyto pozice. Začnu s objektem T, který je už blízko.} \\
\\ 
\texttt{POHYB: T, (3,1)} \\
            }
        }
    }
    & & \\ \\

    \theutterance \stepcounter{utterance}  
    & & \multicolumn{4}{p{0.6\linewidth}}{
        \cellcolor[rgb]{0.9,0.9,0.9}{
            \makecell[{{p{\linewidth}}}]{
                \texttt{\tiny{[P1$\langle$GM]}}
                \texttt{Trest: Vaše zpráva obsahuje více než jeden příkaz!} \\
\\ 
\texttt{Ujistěte se, že vaše odpověď obsahuje pouze `ŘEKNI: <ZPRÁVA>` nebo `POHYB: <OBJEKT>, (<X>, <Y>)` a nic jiného!} \\
\\ 
\texttt{Dohromady jste nasbírali 1 z 12 pokut.} \\
\texttt{Zkuste to prosím znovu!} \\
            }
        }
    }
    & \\ \\

    \theutterance \stepcounter{utterance}  
    & \multicolumn{4}{p{0.6\linewidth}}{
        \cellcolor[rgb]{0.9,0.9,0.9}{
            \makecell[{{p{\linewidth}}}]{
                \texttt{\tiny{[P1$\rangle$GM]}}
                \texttt{POHYB: T, (3,1)} \\
            }
        }
    }
    & & \\ \\

    \theutterance \stepcounter{utterance}  
    & & & \multicolumn{2}{p{0.3\linewidth}}{
        \cellcolor[rgb]{0.9,0.9,0.9}{
            \makecell[{{p{\linewidth}}}]{
                \texttt{\tiny{[GM$|$GM]}}
                \texttt{POHYB: T, (3,1)} \\
            }
        }
    }
    & & \\ \\

    \theutterance \stepcounter{utterance}  
    & & & \multicolumn{2}{p{0.3\linewidth}}{
        \cellcolor[rgb]{0.9,0.9,0.9}{
            \makecell[{{p{\linewidth}}}]{
                \texttt{\tiny{[GM$|$GM]}}
                \texttt{Penalizace: (3,1) není prázdný, ale obsahuje '│'.} \\
\\ 
\texttt{Dohromady jste nasbírali 2 z 12 pokut.} \\
\texttt{Zkuste to prosím znovu!} \\
            }
        }
    }
    & & \\ \\

    \theutterance \stepcounter{utterance}  
    & & \multicolumn{4}{p{0.6\linewidth}}{
        \cellcolor[rgb]{0.9,0.9,0.9}{
            \makecell[{{p{\linewidth}}}]{
                \texttt{\tiny{[P1$\langle$GM]}}
                \texttt{Penalizace: (3,1) není prázdný, ale obsahuje '│'.} \\
\\ 
\texttt{Dohromady jste nasbírali 2 z 12 pokut.} \\
\texttt{Zkuste to prosím znovu!} \\
            }
        }
    }
    & \\ \\

    \theutterance \stepcounter{utterance}  
    & \multicolumn{4}{p{0.6\linewidth}}{
        \cellcolor[rgb]{0.9,0.9,0.9}{
            \makecell[{{p{\linewidth}}}]{
                \texttt{\tiny{[P1$\rangle$GM]}}
                \texttt{POHYB: T, (4,2)} \\
            }
        }
    }
    & & \\ \\

    \theutterance \stepcounter{utterance}  
    & & & \multicolumn{2}{p{0.3\linewidth}}{
        \cellcolor[rgb]{0.9,0.9,0.9}{
            \makecell[{{p{\linewidth}}}]{
                \texttt{\tiny{[GM$|$GM]}}
                \texttt{POHYB: T, (4,2)} \\
            }
        }
    }
    & & \\ \\

    \theutterance \stepcounter{utterance}  
    & & & \multicolumn{2}{p{0.3\linewidth}}{
        \cellcolor[rgb]{0.9,0.9,0.9}{
            \makecell[{{p{\linewidth}}}]{
                \texttt{\tiny{[GM$|$GM]}}
                \texttt{Penalizace: (4,2) není prázdný, ale obsahuje '─'.} \\
\\ 
\texttt{Dohromady jste nasbírali 3 z 12 pokut.} \\
\texttt{Zkuste to prosím znovu!} \\
            }
        }
    }
    & & \\ \\

    \theutterance \stepcounter{utterance}  
    & & \multicolumn{4}{p{0.6\linewidth}}{
        \cellcolor[rgb]{0.9,0.9,0.9}{
            \makecell[{{p{\linewidth}}}]{
                \texttt{\tiny{[P1$\langle$GM]}}
                \texttt{Penalizace: (4,2) není prázdný, ale obsahuje '─'.} \\
\\ 
\texttt{Dohromady jste nasbírali 3 z 12 pokut.} \\
\texttt{Zkuste to prosím znovu!} \\
            }
        }
    }
    & \\ \\

    \theutterance \stepcounter{utterance}  
    & \multicolumn{4}{p{0.6\linewidth}}{
        \cellcolor[rgb]{0.9,0.9,0.9}{
            \makecell[{{p{\linewidth}}}]{
                \texttt{\tiny{[P1$\rangle$GM]}}
                \texttt{POHYB: T, (5,1)} \\
            }
        }
    }
    & & \\ \\

    \theutterance \stepcounter{utterance}  
    & & & \multicolumn{2}{p{0.3\linewidth}}{
        \cellcolor[rgb]{0.9,0.9,0.9}{
            \makecell[{{p{\linewidth}}}]{
                \texttt{\tiny{[GM$|$GM]}}
                \texttt{POHYB: T, (5,1)} \\
            }
        }
    }
    & & \\ \\

    \theutterance \stepcounter{utterance}  
    & & & \multicolumn{2}{p{0.3\linewidth}}{
        \cellcolor[rgb]{0.9,0.9,0.9}{
            \makecell[{{p{\linewidth}}}]{
                \texttt{\tiny{[GM$|$GM]}}
                \texttt{T moved to (5, 1)} \\
            }
        }
    }
    & & \\ \\

    \theutterance \stepcounter{utterance}  
    & & \multicolumn{4}{p{0.6\linewidth}}{
        \cellcolor[rgb]{0.9,0.9,0.9}{
            \makecell[{{p{\linewidth}}}]{
                \texttt{\tiny{[GM$\rangle$P2]}}
                \texttt{Vaše zpráva byla předána druhému hráči.} \\
\\ 
\\ 
\texttt{Právě hrajete kolo 2 z maximálně 20.} \\
\texttt{Dohromady jste nasbírali 3 z 12 pokut.} \\
\texttt{Druhý hráč přesunul objekt na své mřížce.} \\
\\ 
\\ 
\texttt{Jaký je váš další příkaz?} \\
            }
        }
    }
    & \\ \\

    \theutterance \stepcounter{utterance}  
    & & & \multicolumn{4}{p{0.6\linewidth}}{
        \cellcolor[rgb]{0.9,0.9,0.9}{
            \makecell[{{p{\linewidth}}}]{
                \texttt{\tiny{[GM$\langle$P2]}}
                \texttt{Abych mohl pokračovat v plánování, musím vědět, jakou reakci druhý hráč měl na můj návrh a jaký pohyb provedl. Protože mi nebyla sdělena jeho zpráva ani to, jaký objekt přesunul, musím se nejprve zeptat na jeho stav.} \\
\\ 
\texttt{ŘEKNI: Rozumím. Jaký objekt jsi přesunul a jaká je tvá aktuální konfigurace? Také, souhlasíš s mou navrženou cílovou konfigurací W(1,1), I(2,1), T(3,1), C(4,1), H(5,1)?} \\
            }
        }
    }
     \\ \\

    \theutterance \stepcounter{utterance}  
    & & & \multicolumn{2}{p{0.3\linewidth}}{
        \cellcolor[rgb]{0.9,0.9,0.9}{
            \makecell[{{p{\linewidth}}}]{
                \texttt{\tiny{[GM$|$GM]}}
                \texttt{Abych mohl pokračovat v plánování, musím vědět, jakou reakci druhý hráč měl na můj návrh a jaký pohyb provedl. Protože mi nebyla sdělena jeho zpráva ani to, jaký objekt přesunul, musím se nejprve zeptat na jeho stav.} \\
\\ 
\texttt{ŘEKNI: Rozumím. Jaký objekt jsi přesunul a jaká je tvá aktuální konfigurace? Také, souhlasíš s mou navrženou cílovou konfigurací W(1,1), I(2,1), T(3,1), C(4,1), H(5,1)?} \\
            }
        }
    }
    & & \\ \\

    \theutterance \stepcounter{utterance}  
    & & & \multicolumn{2}{p{0.3\linewidth}}{
        \cellcolor[rgb]{0.9,0.9,0.9}{
            \makecell[{{p{\linewidth}}}]{
                \texttt{\tiny{[GM$|$GM]}}
                \texttt{Invalid format: head is not empty: 'Abych mohl pokračovat v plánování, musím vědět, jakou reakci druhý hráč měl na můj návrh a jaký pohyb provedl. Protože mi nebyla sdělena jeho zpráva ani to, jaký objekt přesunul, musím se nejprve zeptat na jeho stav.} \\
\\ 
\texttt{'} \\
            }
        }
    }
    & & \\ \\

    \theutterance \stepcounter{utterance}  
    & & \multicolumn{4}{p{0.6\linewidth}}{
        \cellcolor[rgb]{0.9,0.9,0.9}{
            \makecell[{{p{\linewidth}}}]{
                \texttt{\tiny{[GM$\rangle$P2]}}
                \texttt{Trest: Vaše zpráva nesmí obsahovat nic před příkazem!} \\
\\ 
\texttt{Ujistěte se, že vaše odpověď obsahuje pouze `ŘEKNI: <ZPRÁVA>` nebo `POHYB: <OBJEKT>, (<X>, <Y>)` a nic jiného!} \\
\\ 
\texttt{Dohromady jste nasbírali 4 z 12 pokut.} \\
\texttt{Zkuste to prosím znovu!} \\
            }
        }
    }
    & \\ \\

    \theutterance \stepcounter{utterance}  
    & & & \multicolumn{4}{p{0.6\linewidth}}{
        \cellcolor[rgb]{0.9,0.9,0.9}{
            \makecell[{{p{\linewidth}}}]{
                \texttt{\tiny{[GM$\langle$P2]}}
                \texttt{ŘEKNI: Rozumím. Jaký objekt jsi přesunul a jaká je tvá aktuální konfigurace? Také, souhlasíš s mou navrženou cílovou konfigurací W(1,1), I(2,1), T(3,1), C(4,1), H(5,1)?} \\
            }
        }
    }
     \\ \\

    \theutterance \stepcounter{utterance}  
    & & & \multicolumn{2}{p{0.3\linewidth}}{
        \cellcolor[rgb]{0.9,0.9,0.9}{
            \makecell[{{p{\linewidth}}}]{
                \texttt{\tiny{[GM$|$GM]}}
                \texttt{ŘEKNI: Rozumím. Jaký objekt jsi přesunul a jaká je tvá aktuální konfigurace? Také, souhlasíš s mou navrženou cílovou konfigurací W(1,1), I(2,1), T(3,1), C(4,1), H(5,1)?} \\
            }
        }
    }
    & & \\ \\

    \theutterance \stepcounter{utterance}  
    & & \multicolumn{4}{p{0.6\linewidth}}{
        \cellcolor[rgb]{0.9,0.9,0.9}{
            \makecell[{{p{\linewidth}}}]{
                \texttt{\tiny{[P1$\langle$GM]}}
                \texttt{Úspěšně přesunuto 'T' do (5,1). Vaše aktualizovaná mřížka vypadá takto:} \\
\texttt{```} \\
\texttt{ 1234567} \\
\texttt{╔══╤══╤═╗} \\
\texttt{║◌◌│◌T│◌║ 1} \\
\texttt{╟──┼──┴─╢ 2} \\
\texttt{║WI│◌◌◌◌║ 3} \\
\texttt{╟──┘◌◌◌◌║ 4} \\
\texttt{║◌◌◌H◌◌C║ 5} \\
\texttt{║◌◌┌─┬──╢ 6} \\
\texttt{║◌◌│◌│◌◌║ 7} \\
\texttt{╚══╧═╧══╝} \\
\texttt{```} \\
\texttt{Právě hrajete kolo 2 z maximálně 20.} \\
\texttt{Dohromady jste nasbírali 4 z 12 pokut.} \\
\texttt{Druhý hráč poslal tuto zprávu:} \\
\\ 
\\ 
\texttt{"Rozumím. Jaký objekt jsi přesunul a jaká je tvá aktuální konfigurace? Také, souhlasíš s mou navrženou cílovou konfigurací W(1,1), I(2,1), T(3,1), C(4,1), H(5,1)?"} \\
\\ 
\\ 
\texttt{Jaký je váš další příkaz?} \\
            }
        }
    }
    & \\ \\

    \theutterance \stepcounter{utterance}  
    & \multicolumn{4}{p{0.6\linewidth}}{
        \cellcolor[rgb]{0.9,0.9,0.9}{
            \makecell[{{p{\linewidth}}}]{
                \texttt{\tiny{[P1$\rangle$GM]}}
                \texttt{ŘEKNI: Přesunul jsem T na (5,1). Moje aktuální pozice jsou W(1,3), I(2,3), T(5,1), H(4,5), C(7,5). Souhlasím s tvým cílem W(1,1), I(2,1), T(3,1), C(4,1), H(5,1), ale musíme si uvědomit, že T je teď na (5,1), což je tvůj cíl pro H. Budeme muset T později posunout na (3,1) a H na (5,1). Pokračuj, já se budu snažit o své objekty.} \\
            }
        }
    }
    & & \\ \\

    \theutterance \stepcounter{utterance}  
    & & & \multicolumn{2}{p{0.3\linewidth}}{
        \cellcolor[rgb]{0.9,0.9,0.9}{
            \makecell[{{p{\linewidth}}}]{
                \texttt{\tiny{[GM$|$GM]}}
                \texttt{ŘEKNI: Přesunul jsem T na (5,1). Moje aktuální pozice jsou W(1,3), I(2,3), T(5,1), H(4,5), C(7,5). Souhlasím s tvým cílem W(1,1), I(2,1), T(3,1), C(4,1), H(5,1), ale musíme si uvědomit, že T je teď na (5,1), což je tvůj cíl pro H. Budeme muset T později posunout na (3,1) a H na (5,1). Pokračuj, já se budu snažit o své objekty.} \\
            }
        }
    }
    & & \\ \\

    \theutterance \stepcounter{utterance}  
    & & \multicolumn{4}{p{0.6\linewidth}}{
        \cellcolor[rgb]{0.9,0.9,0.9}{
            \makecell[{{p{\linewidth}}}]{
                \texttt{\tiny{[GM$\rangle$P2]}}
                \texttt{Vaše zpráva byla předána druhému hráči.} \\
\\ 
\\ 
\texttt{Právě hrajete kolo 3 z maximálně 20.} \\
\texttt{Dohromady jste nasbírali 4 z 12 pokut.} \\
\texttt{Druhý hráč poslal tuto zprávu:} \\
\\ 
\\ 
\texttt{"Přesunul jsem T na (5,1). Moje aktuální pozice jsou W(1,3), I(2,3), T(5,1), H(4,5), C(7,5). Souhlasím s tvým cílem W(1,1), I(2,1), T(3,1), C(4,1), H(5,1), ale musíme si uvědomit, že T je teď na (5,1), což je tvůj cíl pro H. Budeme muset T později posunout na (3,1) a H na (5,1). Pokračuj, já se budu snažit o své objekty."} \\
\\ 
\\ 
\texttt{Jaký je váš další příkaz?} \\
            }
        }
    }
    & \\ \\

    \theutterance \stepcounter{utterance}  
    & & & \multicolumn{4}{p{0.6\linewidth}}{
        \cellcolor[rgb]{0.9,0.9,0.9}{
            \makecell[{{p{\linewidth}}}]{
                \texttt{\tiny{[GM$\langle$P2]}}
                \texttt{POHYB: W, (1, 1)} \\
            }
        }
    }
     \\ \\

    \theutterance \stepcounter{utterance}  
    & & & \multicolumn{2}{p{0.3\linewidth}}{
        \cellcolor[rgb]{0.9,0.9,0.9}{
            \makecell[{{p{\linewidth}}}]{
                \texttt{\tiny{[GM$|$GM]}}
                \texttt{POHYB: W, (1, 1)} \\
            }
        }
    }
    & & \\ \\

    \theutterance \stepcounter{utterance}  
    & & & \multicolumn{2}{p{0.3\linewidth}}{
        \cellcolor[rgb]{0.9,0.9,0.9}{
            \makecell[{{p{\linewidth}}}]{
                \texttt{\tiny{[GM$|$GM]}}
                \texttt{W moved to (1, 1)} \\
            }
        }
    }
    & & \\ \\

    \theutterance \stepcounter{utterance}  
    & & \multicolumn{4}{p{0.6\linewidth}}{
        \cellcolor[rgb]{0.9,0.9,0.9}{
            \makecell[{{p{\linewidth}}}]{
                \texttt{\tiny{[P1$\langle$GM]}}
                \texttt{Vaše zpráva byla předána druhému hráči.} \\
\\ 
\\ 
\texttt{Právě hrajete kolo 3 z maximálně 20.} \\
\texttt{Dohromady jste nasbírali 4 z 12 pokut.} \\
\texttt{Druhý hráč přesunul objekt na své mřížce.} \\
\\ 
\\ 
\texttt{Jaký je váš další příkaz?} \\
            }
        }
    }
    & \\ \\

    \theutterance \stepcounter{utterance}  
    & \multicolumn{4}{p{0.6\linewidth}}{
        \cellcolor[rgb]{0.9,0.9,0.9}{
            \makecell[{{p{\linewidth}}}]{
                \texttt{\tiny{[P1$\rangle$GM]}}
                \texttt{ŘEKNI: Jaký pohyb jsi udělal? Já se nyní pokusím posunout W na (1,1).} \\
            }
        }
    }
    & & \\ \\

    \theutterance \stepcounter{utterance}  
    & & & \multicolumn{2}{p{0.3\linewidth}}{
        \cellcolor[rgb]{0.9,0.9,0.9}{
            \makecell[{{p{\linewidth}}}]{
                \texttt{\tiny{[GM$|$GM]}}
                \texttt{ŘEKNI: Jaký pohyb jsi udělal? Já se nyní pokusím posunout W na (1,1).} \\
            }
        }
    }
    & & \\ \\

    \theutterance \stepcounter{utterance}  
    & & \multicolumn{4}{p{0.6\linewidth}}{
        \cellcolor[rgb]{0.9,0.9,0.9}{
            \makecell[{{p{\linewidth}}}]{
                \texttt{\tiny{[GM$\rangle$P2]}}
                \texttt{Úspěšně přesunuto 'W' do (1,1). Vaše aktualizovaná mřížka vypadá takto:} \\
\texttt{```} \\
\texttt{ 1234567} \\
\texttt{╔══╤══╤═╗} \\
\texttt{║W◌│◌◌│I║ 1} \\
\texttt{╟──┼──┴─╢ 2} \\
\texttt{║◌◌│◌◌◌◌║ 3} \\
\texttt{╟──┘T◌◌◌║ 4} \\
\texttt{║H◌◌◌◌◌◌║ 5} \\
\texttt{║◌◌┌─┬──╢ 6} \\
\texttt{║C◌│◌│◌◌║ 7} \\
\texttt{╚══╧═╧══╝} \\
\texttt{```} \\
\texttt{Právě hrajete kolo 4 z maximálně 20.} \\
\texttt{Dohromady jste nasbírali 4 z 12 pokut.} \\
\texttt{Druhý hráč poslal tuto zprávu:} \\
\\ 
\\ 
\texttt{"Jaký pohyb jsi udělal? Já se nyní pokusím posunout W na (1,1)."} \\
\\ 
\\ 
\texttt{Jaký je váš další příkaz?} \\
            }
        }
    }
    & \\ \\

    \theutterance \stepcounter{utterance}  
    & & & \multicolumn{4}{p{0.6\linewidth}}{
        \cellcolor[rgb]{0.9,0.9,0.9}{
            \makecell[{{p{\linewidth}}}]{
                \texttt{\tiny{[GM$\langle$P2]}}
                \texttt{ŘEKNI: Udělal jsem pohyb W na (1,1). Teď se pokusím posunout I na (2,1).} \\
            }
        }
    }
     \\ \\

    \theutterance \stepcounter{utterance}  
    & & & \multicolumn{2}{p{0.3\linewidth}}{
        \cellcolor[rgb]{0.9,0.9,0.9}{
            \makecell[{{p{\linewidth}}}]{
                \texttt{\tiny{[GM$|$GM]}}
                \texttt{ŘEKNI: Udělal jsem pohyb W na (1,1). Teď se pokusím posunout I na (2,1).} \\
            }
        }
    }
    & & \\ \\

    \theutterance \stepcounter{utterance}  
    & & \multicolumn{4}{p{0.6\linewidth}}{
        \cellcolor[rgb]{0.9,0.9,0.9}{
            \makecell[{{p{\linewidth}}}]{
                \texttt{\tiny{[P1$\langle$GM]}}
                \texttt{Vaše zpráva byla předána druhému hráči.} \\
\\ 
\\ 
\texttt{Právě hrajete kolo 4 z maximálně 20.} \\
\texttt{Dohromady jste nasbírali 4 z 12 pokut.} \\
\texttt{Druhý hráč poslal tuto zprávu:} \\
\\ 
\\ 
\texttt{"Udělal jsem pohyb W na (1,1). Teď se pokusím posunout I na (2,1)."} \\
\\ 
\\ 
\texttt{Jaký je váš další příkaz?} \\
            }
        }
    }
    & \\ \\

    \theutterance \stepcounter{utterance}  
    & \multicolumn{4}{p{0.6\linewidth}}{
        \cellcolor[rgb]{0.9,0.9,0.9}{
            \makecell[{{p{\linewidth}}}]{
                \texttt{\tiny{[P1$\rangle$GM]}}
                \texttt{POHYB: W, (1,1)} \\
            }
        }
    }
    & & \\ \\

    \theutterance \stepcounter{utterance}  
    & & & \multicolumn{2}{p{0.3\linewidth}}{
        \cellcolor[rgb]{0.9,0.9,0.9}{
            \makecell[{{p{\linewidth}}}]{
                \texttt{\tiny{[GM$|$GM]}}
                \texttt{POHYB: W, (1,1)} \\
            }
        }
    }
    & & \\ \\

    \theutterance \stepcounter{utterance}  
    & & & \multicolumn{2}{p{0.3\linewidth}}{
        \cellcolor[rgb]{0.9,0.9,0.9}{
            \makecell[{{p{\linewidth}}}]{
                \texttt{\tiny{[GM$|$GM]}}
                \texttt{W moved to (1, 1)} \\
            }
        }
    }
    & & \\ \\

    \theutterance \stepcounter{utterance}  
    & & \multicolumn{4}{p{0.6\linewidth}}{
        \cellcolor[rgb]{0.9,0.9,0.9}{
            \makecell[{{p{\linewidth}}}]{
                \texttt{\tiny{[GM$\rangle$P2]}}
                \texttt{Vaše zpráva byla předána druhému hráči.} \\
\\ 
\\ 
\texttt{Právě hrajete kolo 5 z maximálně 20.} \\
\texttt{Dohromady jste nasbírali 4 z 12 pokut.} \\
\texttt{Druhý hráč přesunul objekt na své mřížce.} \\
\\ 
\\ 
\texttt{Jaký je váš další příkaz?} \\
            }
        }
    }
    & \\ \\

    \theutterance \stepcounter{utterance}  
    & & & \multicolumn{4}{p{0.6\linewidth}}{
        \cellcolor[rgb]{0.9,0.9,0.9}{
            \makecell[{{p{\linewidth}}}]{
                \texttt{\tiny{[GM$\langle$P2]}}
                \texttt{ŘEKNI: Jaký pohyb jsi udělal? Moje aktuální pozice jsou W(1,1), I(7,1), T(4,4), C(1,7), H(1,5).} \\
            }
        }
    }
     \\ \\

    \theutterance \stepcounter{utterance}  
    & & & \multicolumn{2}{p{0.3\linewidth}}{
        \cellcolor[rgb]{0.9,0.9,0.9}{
            \makecell[{{p{\linewidth}}}]{
                \texttt{\tiny{[GM$|$GM]}}
                \texttt{ŘEKNI: Jaký pohyb jsi udělal? Moje aktuální pozice jsou W(1,1), I(7,1), T(4,4), C(1,7), H(1,5).} \\
            }
        }
    }
    & & \\ \\

    \theutterance \stepcounter{utterance}  
    & & \multicolumn{4}{p{0.6\linewidth}}{
        \cellcolor[rgb]{0.9,0.9,0.9}{
            \makecell[{{p{\linewidth}}}]{
                \texttt{\tiny{[P1$\langle$GM]}}
                \texttt{Úspěšně přesunuto 'W' do (1,1). Vaše aktualizovaná mřížka vypadá takto:} \\
\texttt{```} \\
\texttt{ 1234567} \\
\texttt{╔══╤══╤═╗} \\
\texttt{║W◌│◌T│◌║ 1} \\
\texttt{╟──┼──┴─╢ 2} \\
\texttt{║◌I│◌◌◌◌║ 3} \\
\texttt{╟──┘◌◌◌◌║ 4} \\
\texttt{║◌◌◌H◌◌C║ 5} \\
\texttt{║◌◌┌─┬──╢ 6} \\
\texttt{║◌◌│◌│◌◌║ 7} \\
\texttt{╚══╧═╧══╝} \\
\texttt{```} \\
\texttt{Právě hrajete kolo 5 z maximálně 20.} \\
\texttt{Dohromady jste nasbírali 4 z 12 pokut.} \\
\texttt{Druhý hráč poslal tuto zprávu:} \\
\\ 
\\ 
\texttt{"Jaký pohyb jsi udělal? Moje aktuální pozice jsou W(1,1), I(7,1), T(4,4), C(1,7), H(1,5)."} \\
\\ 
\\ 
\texttt{Jaký je váš další příkaz?} \\
            }
        }
    }
    & \\ \\

    \theutterance \stepcounter{utterance}  
    & \multicolumn{4}{p{0.6\linewidth}}{
        \cellcolor[rgb]{0.9,0.9,0.9}{
            \makecell[{{p{\linewidth}}}]{
                \texttt{\tiny{[P1$\rangle$GM]}}
                \texttt{POHYB: I, (2,1)} \\
            }
        }
    }
    & & \\ \\

    \theutterance \stepcounter{utterance}  
    & & & \multicolumn{2}{p{0.3\linewidth}}{
        \cellcolor[rgb]{0.9,0.9,0.9}{
            \makecell[{{p{\linewidth}}}]{
                \texttt{\tiny{[GM$|$GM]}}
                \texttt{POHYB: I, (2,1)} \\
            }
        }
    }
    & & \\ \\

    \theutterance \stepcounter{utterance}  
    & & & \multicolumn{2}{p{0.3\linewidth}}{
        \cellcolor[rgb]{0.9,0.9,0.9}{
            \makecell[{{p{\linewidth}}}]{
                \texttt{\tiny{[GM$|$GM]}}
                \texttt{I moved to (2, 1)} \\
            }
        }
    }
    & & \\ \\

    \theutterance \stepcounter{utterance}  
    & & \multicolumn{4}{p{0.6\linewidth}}{
        \cellcolor[rgb]{0.9,0.9,0.9}{
            \makecell[{{p{\linewidth}}}]{
                \texttt{\tiny{[GM$\rangle$P2]}}
                \texttt{Vaše zpráva byla předána druhému hráči.} \\
\\ 
\\ 
\texttt{Právě hrajete kolo 6 z maximálně 20.} \\
\texttt{Dohromady jste nasbírali 4 z 12 pokut.} \\
\texttt{Druhý hráč přesunul objekt na své mřížce.} \\
\\ 
\\ 
\texttt{Jaký je váš další příkaz?} \\
            }
        }
    }
    & \\ \\

    \theutterance \stepcounter{utterance}  
    & & & \multicolumn{4}{p{0.6\linewidth}}{
        \cellcolor[rgb]{0.9,0.9,0.9}{
            \makecell[{{p{\linewidth}}}]{
                \texttt{\tiny{[GM$\langle$P2]}}
                \texttt{ŘEKNI: Jaký pohyb jsi udělal? Potřebuji znát tvou aktuální konfiguraci, abychom věděli, jak moc jsme se přiblížili cíli.} \\
            }
        }
    }
     \\ \\

    \theutterance \stepcounter{utterance}  
    & & & \multicolumn{2}{p{0.3\linewidth}}{
        \cellcolor[rgb]{0.9,0.9,0.9}{
            \makecell[{{p{\linewidth}}}]{
                \texttt{\tiny{[GM$|$GM]}}
                \texttt{ŘEKNI: Jaký pohyb jsi udělal? Potřebuji znát tvou aktuální konfiguraci, abychom věděli, jak moc jsme se přiblížili cíli.} \\
            }
        }
    }
    & & \\ \\

    \theutterance \stepcounter{utterance}  
    & & \multicolumn{4}{p{0.6\linewidth}}{
        \cellcolor[rgb]{0.9,0.9,0.9}{
            \makecell[{{p{\linewidth}}}]{
                \texttt{\tiny{[P1$\langle$GM]}}
                \texttt{Úspěšně přesunuto 'I' do (2,1). Vaše aktualizovaná mřížka vypadá takto:} \\
\texttt{```} \\
\texttt{ 1234567} \\
\texttt{╔══╤══╤═╗} \\
\texttt{║WI│◌T│◌║ 1} \\
\texttt{╟──┼──┴─╢ 2} \\
\texttt{║◌◌│◌◌◌◌║ 3} \\
\texttt{╟──┘◌◌◌◌║ 4} \\
\texttt{║◌◌◌H◌◌C║ 5} \\
\texttt{║◌◌┌─┬──╢ 6} \\
\texttt{║◌◌│◌│◌◌║ 7} \\
\texttt{╚══╧═╧══╝} \\
\texttt{```} \\
\texttt{Právě hrajete kolo 6 z maximálně 20.} \\
\texttt{Dohromady jste nasbírali 4 z 12 pokut.} \\
\texttt{Druhý hráč poslal tuto zprávu:} \\
\\ 
\\ 
\texttt{"Jaký pohyb jsi udělal? Potřebuji znát tvou aktuální konfiguraci, abychom věděli, jak moc jsme se přiblížili cíli."} \\
\\ 
\\ 
\texttt{Jaký je váš další příkaz?} \\
            }
        }
    }
    & \\ \\

    \theutterance \stepcounter{utterance}  
    & \multicolumn{4}{p{0.6\linewidth}}{
        \cellcolor[rgb]{0.9,0.9,0.9}{
            \makecell[{{p{\linewidth}}}]{
                \texttt{\tiny{[P1$\rangle$GM]}}
                \texttt{ŘEKNI: Přesunul jsem I na (2,1). Moje aktuální pozice jsou W(1,1), I(2,1), T(5,1), H(4,5), C(7,5).} \\
            }
        }
    }
    & & \\ \\

    \theutterance \stepcounter{utterance}  
    & & & \multicolumn{2}{p{0.3\linewidth}}{
        \cellcolor[rgb]{0.9,0.9,0.9}{
            \makecell[{{p{\linewidth}}}]{
                \texttt{\tiny{[GM$|$GM]}}
                \texttt{ŘEKNI: Přesunul jsem I na (2,1). Moje aktuální pozice jsou W(1,1), I(2,1), T(5,1), H(4,5), C(7,5).} \\
            }
        }
    }
    & & \\ \\

    \theutterance \stepcounter{utterance}  
    & & \multicolumn{4}{p{0.6\linewidth}}{
        \cellcolor[rgb]{0.9,0.9,0.9}{
            \makecell[{{p{\linewidth}}}]{
                \texttt{\tiny{[GM$\rangle$P2]}}
                \texttt{Vaše zpráva byla předána druhému hráči.} \\
\\ 
\\ 
\texttt{Právě hrajete kolo 7 z maximálně 20.} \\
\texttt{Dohromady jste nasbírali 4 z 12 pokut.} \\
\texttt{Druhý hráč poslal tuto zprávu:} \\
\\ 
\\ 
\texttt{"Přesunul jsem I na (2,1). Moje aktuální pozice jsou W(1,1), I(2,1), T(5,1), H(4,5), C(7,5)."} \\
\\ 
\\ 
\texttt{Jaký je váš další příkaz?} \\
            }
        }
    }
    & \\ \\

    \theutterance \stepcounter{utterance}  
    & & & \multicolumn{4}{p{0.6\linewidth}}{
        \cellcolor[rgb]{0.9,0.9,0.9}{
            \makecell[{{p{\linewidth}}}]{
                \texttt{\tiny{[GM$\langle$P2]}}
                \texttt{POHYB: I, (2, 1)} \\
            }
        }
    }
     \\ \\

    \theutterance \stepcounter{utterance}  
    & & & \multicolumn{2}{p{0.3\linewidth}}{
        \cellcolor[rgb]{0.9,0.9,0.9}{
            \makecell[{{p{\linewidth}}}]{
                \texttt{\tiny{[GM$|$GM]}}
                \texttt{POHYB: I, (2, 1)} \\
            }
        }
    }
    & & \\ \\

    \theutterance \stepcounter{utterance}  
    & & & \multicolumn{2}{p{0.3\linewidth}}{
        \cellcolor[rgb]{0.9,0.9,0.9}{
            \makecell[{{p{\linewidth}}}]{
                \texttt{\tiny{[GM$|$GM]}}
                \texttt{I moved to (2, 1)} \\
            }
        }
    }
    & & \\ \\

    \theutterance \stepcounter{utterance}  
    & & \multicolumn{4}{p{0.6\linewidth}}{
        \cellcolor[rgb]{0.9,0.9,0.9}{
            \makecell[{{p{\linewidth}}}]{
                \texttt{\tiny{[P1$\langle$GM]}}
                \texttt{Vaše zpráva byla předána druhému hráči.} \\
\\ 
\\ 
\texttt{Právě hrajete kolo 7 z maximálně 20.} \\
\texttt{Dohromady jste nasbírali 4 z 12 pokut.} \\
\texttt{Druhý hráč přesunul objekt na své mřížce.} \\
\\ 
\\ 
\texttt{Jaký je váš další příkaz?} \\
            }
        }
    }
    & \\ \\

    \theutterance \stepcounter{utterance}  
    & \multicolumn{4}{p{0.6\linewidth}}{
        \cellcolor[rgb]{0.9,0.9,0.9}{
            \makecell[{{p{\linewidth}}}]{
                \texttt{\tiny{[P1$\rangle$GM]}}
                \texttt{ŘEKNI: Jaký pohyb jsi udělal? Moje aktuální pozice jsou W(1,1), I(2,1), T(5,1), H(4,5), C(7,5). Teď se pokusím posunout H na (5,1), ale musím nejdřív uvolnit místo pro T. Takže nejprve posuneme T na (3,1).} \\
            }
        }
    }
    & & \\ \\

    \theutterance \stepcounter{utterance}  
    & & & \multicolumn{2}{p{0.3\linewidth}}{
        \cellcolor[rgb]{0.9,0.9,0.9}{
            \makecell[{{p{\linewidth}}}]{
                \texttt{\tiny{[GM$|$GM]}}
                \texttt{ŘEKNI: Jaký pohyb jsi udělal? Moje aktuální pozice jsou W(1,1), I(2,1), T(5,1), H(4,5), C(7,5). Teď se pokusím posunout H na (5,1), ale musím nejdřív uvolnit místo pro T. Takže nejprve posuneme T na (3,1).} \\
            }
        }
    }
    & & \\ \\

    \theutterance \stepcounter{utterance}  
    & & \multicolumn{4}{p{0.6\linewidth}}{
        \cellcolor[rgb]{0.9,0.9,0.9}{
            \makecell[{{p{\linewidth}}}]{
                \texttt{\tiny{[GM$\rangle$P2]}}
                \texttt{Úspěšně přesunuto 'I' do (2,1). Vaše aktualizovaná mřížka vypadá takto:} \\
\texttt{```} \\
\texttt{ 1234567} \\
\texttt{╔══╤══╤═╗} \\
\texttt{║WI│◌◌│◌║ 1} \\
\texttt{╟──┼──┴─╢ 2} \\
\texttt{║◌◌│◌◌◌◌║ 3} \\
\texttt{╟──┘T◌◌◌║ 4} \\
\texttt{║H◌◌◌◌◌◌║ 5} \\
\texttt{║◌◌┌─┬──╢ 6} \\
\texttt{║C◌│◌│◌◌║ 7} \\
\texttt{╚══╧═╧══╝} \\
\texttt{```} \\
\texttt{Právě hrajete kolo 8 z maximálně 20.} \\
\texttt{Dohromady jste nasbírali 4 z 12 pokut.} \\
\texttt{Druhý hráč poslal tuto zprávu:} \\
\\ 
\\ 
\texttt{"Jaký pohyb jsi udělal? Moje aktuální pozice jsou W(1,1), I(2,1), T(5,1), H(4,5), C(7,5). Teď se pokusím posunout H na (5,1), ale musím nejdřív uvolnit místo pro T. Takže nejprve posuneme T na (3,1)."} \\
\\ 
\\ 
\texttt{Jaký je váš další příkaz?} \\
            }
        }
    }
    & \\ \\

    \theutterance \stepcounter{utterance}  
    & & & \multicolumn{4}{p{0.6\linewidth}}{
        \cellcolor[rgb]{0.9,0.9,0.9}{
            \makecell[{{p{\linewidth}}}]{
                \texttt{\tiny{[GM$\langle$P2]}}
                \texttt{ŘEKNI: Udělal jsem pohyb I na (2,1). To je skvělé, máme W i I na správných místech! Já se teď pokusím posunout T na (3,1).} \\
            }
        }
    }
     \\ \\

    \theutterance \stepcounter{utterance}  
    & & & \multicolumn{2}{p{0.3\linewidth}}{
        \cellcolor[rgb]{0.9,0.9,0.9}{
            \makecell[{{p{\linewidth}}}]{
                \texttt{\tiny{[GM$|$GM]}}
                \texttt{ŘEKNI: Udělal jsem pohyb I na (2,1). To je skvělé, máme W i I na správných místech! Já se teď pokusím posunout T na (3,1).} \\
            }
        }
    }
    & & \\ \\

    \theutterance \stepcounter{utterance}  
    & & \multicolumn{4}{p{0.6\linewidth}}{
        \cellcolor[rgb]{0.9,0.9,0.9}{
            \makecell[{{p{\linewidth}}}]{
                \texttt{\tiny{[P1$\langle$GM]}}
                \texttt{Vaše zpráva byla předána druhému hráči.} \\
\\ 
\\ 
\texttt{Právě hrajete kolo 8 z maximálně 20.} \\
\texttt{Dohromady jste nasbírali 4 z 12 pokut.} \\
\texttt{Druhý hráč poslal tuto zprávu:} \\
\\ 
\\ 
\texttt{"Udělal jsem pohyb I na (2,1). To je skvělé, máme W i I na správných místech! Já se teď pokusím posunout T na (3,1)."} \\
\\ 
\\ 
\texttt{Jaký je váš další příkaz?} \\
            }
        }
    }
    & \\ \\

    \theutterance \stepcounter{utterance}  
    & \multicolumn{4}{p{0.6\linewidth}}{
        \cellcolor[rgb]{0.9,0.9,0.9}{
            \makecell[{{p{\linewidth}}}]{
                \texttt{\tiny{[P1$\rangle$GM]}}
                \texttt{POHYB: T, (3,1)} \\
            }
        }
    }
    & & \\ \\

    \theutterance \stepcounter{utterance}  
    & & & \multicolumn{2}{p{0.3\linewidth}}{
        \cellcolor[rgb]{0.9,0.9,0.9}{
            \makecell[{{p{\linewidth}}}]{
                \texttt{\tiny{[GM$|$GM]}}
                \texttt{POHYB: T, (3,1)} \\
            }
        }
    }
    & & \\ \\

    \theutterance \stepcounter{utterance}  
    & & & \multicolumn{2}{p{0.3\linewidth}}{
        \cellcolor[rgb]{0.9,0.9,0.9}{
            \makecell[{{p{\linewidth}}}]{
                \texttt{\tiny{[GM$|$GM]}}
                \texttt{Penalizace: (3,1) není prázdný, ale obsahuje '│'.} \\
\\ 
\texttt{Dohromady jste nasbírali 5 z 12 pokut.} \\
\texttt{Zkuste to prosím znovu!} \\
            }
        }
    }
    & & \\ \\

    \theutterance \stepcounter{utterance}  
    & & \multicolumn{4}{p{0.6\linewidth}}{
        \cellcolor[rgb]{0.9,0.9,0.9}{
            \makecell[{{p{\linewidth}}}]{
                \texttt{\tiny{[P1$\langle$GM]}}
                \texttt{Penalizace: (3,1) není prázdný, ale obsahuje '│'.} \\
\\ 
\texttt{Dohromady jste nasbírali 5 z 12 pokut.} \\
\texttt{Zkuste to prosím znovu!} \\
            }
        }
    }
    & \\ \\

    \theutterance \stepcounter{utterance}  
    & \multicolumn{4}{p{0.6\linewidth}}{
        \cellcolor[rgb]{0.9,0.9,0.9}{
            \makecell[{{p{\linewidth}}}]{
                \texttt{\tiny{[P1$\rangle$GM]}}
                \texttt{POHYB: T, (6,1)} \\
            }
        }
    }
    & & \\ \\

    \theutterance \stepcounter{utterance}  
    & & & \multicolumn{2}{p{0.3\linewidth}}{
        \cellcolor[rgb]{0.9,0.9,0.9}{
            \makecell[{{p{\linewidth}}}]{
                \texttt{\tiny{[GM$|$GM]}}
                \texttt{POHYB: T, (6,1)} \\
            }
        }
    }
    & & \\ \\

    \theutterance \stepcounter{utterance}  
    & & & \multicolumn{2}{p{0.3\linewidth}}{
        \cellcolor[rgb]{0.9,0.9,0.9}{
            \makecell[{{p{\linewidth}}}]{
                \texttt{\tiny{[GM$|$GM]}}
                \texttt{Penalizace: (6,1) není prázdný, ale obsahuje '│'.} \\
\\ 
\texttt{Dohromady jste nasbírali 6 z 12 pokut.} \\
\texttt{Zkuste to prosím znovu!} \\
            }
        }
    }
    & & \\ \\

    \theutterance \stepcounter{utterance}  
    & & \multicolumn{4}{p{0.6\linewidth}}{
        \cellcolor[rgb]{0.9,0.9,0.9}{
            \makecell[{{p{\linewidth}}}]{
                \texttt{\tiny{[P1$\langle$GM]}}
                \texttt{Penalizace: (6,1) není prázdný, ale obsahuje '│'.} \\
\\ 
\texttt{Dohromady jste nasbírali 6 z 12 pokut.} \\
\texttt{Zkuste to prosím znovu!} \\
            }
        }
    }
    & \\ \\

    \theutterance \stepcounter{utterance}  
    & \multicolumn{4}{p{0.6\linewidth}}{
        \cellcolor[rgb]{0.9,0.9,0.9}{
            \makecell[{{p{\linewidth}}}]{
                \texttt{\tiny{[P1$\rangle$GM]}}
                \texttt{POHYB: T, (5,2)} \\
            }
        }
    }
    & & \\ \\

    \theutterance \stepcounter{utterance}  
    & & & \multicolumn{2}{p{0.3\linewidth}}{
        \cellcolor[rgb]{0.9,0.9,0.9}{
            \makecell[{{p{\linewidth}}}]{
                \texttt{\tiny{[GM$|$GM]}}
                \texttt{POHYB: T, (5,2)} \\
            }
        }
    }
    & & \\ \\

    \theutterance \stepcounter{utterance}  
    & & & \multicolumn{2}{p{0.3\linewidth}}{
        \cellcolor[rgb]{0.9,0.9,0.9}{
            \makecell[{{p{\linewidth}}}]{
                \texttt{\tiny{[GM$|$GM]}}
                \texttt{Penalizace: (5,2) není prázdný, ale obsahuje '─'.} \\
\\ 
\texttt{Dohromady jste nasbírali 7 z 12 pokut.} \\
\texttt{Zkuste to prosím znovu!} \\
            }
        }
    }
    & & \\ \\

    \theutterance \stepcounter{utterance}  
    & & \multicolumn{4}{p{0.6\linewidth}}{
        \cellcolor[rgb]{0.9,0.9,0.9}{
            \makecell[{{p{\linewidth}}}]{
                \texttt{\tiny{[P1$\langle$GM]}}
                \texttt{Penalizace: (5,2) není prázdný, ale obsahuje '─'.} \\
\\ 
\texttt{Dohromady jste nasbírali 7 z 12 pokut.} \\
\texttt{Zkuste to prosím znovu!} \\
            }
        }
    }
    & \\ \\

    \theutterance \stepcounter{utterance}  
    & \multicolumn{4}{p{0.6\linewidth}}{
        \cellcolor[rgb]{0.9,0.9,0.9}{
            \makecell[{{p{\linewidth}}}]{
                \texttt{\tiny{[P1$\rangle$GM]}}
                \texttt{POHYB: T, (6,1)} \\
\\ 
\texttt{Wait, I see the grid again. (6,1) is `◌`. Let me re{-}check the grid.} \\
\texttt{The grid at row 1 is: `║W I │ ◌ T │ ◌║ 1`} \\
\texttt{Wait, the columns are:} \\
\texttt{1: W} \\
\texttt{2: I} \\
\texttt{3: │ (boundary)} \\
\texttt{4: ◌} \\
\texttt{5: T} \\
\texttt{6: │ (boundary)} \\
\texttt{7: ◌} \\
\\ 
\texttt{So (3,1) and (6,1) are boundaries. My previous attempts were wrong.} \\
\texttt{Let's try to move T to a safe empty cell.} \\
\texttt{T is at (5,1).} \\
\texttt{(4,1) is `◌`.} \\
\\ 
\texttt{POHYB: T, (4,1)} \\
            }
        }
    }
    & & \\ \\

    \theutterance \stepcounter{utterance}  
    & & & \multicolumn{2}{p{0.3\linewidth}}{
        \cellcolor[rgb]{0.9,0.9,0.9}{
            \makecell[{{p{\linewidth}}}]{
                \texttt{\tiny{[GM$|$GM]}}
                \texttt{POHYB: T, (6,1)} \\
\\ 
\texttt{Wait, I see the grid again. (6,1) is `◌`. Let me re{-}check the grid.} \\
\texttt{The grid at row 1 is: `║W I │ ◌ T │ ◌║ 1`} \\
\texttt{Wait, the columns are:} \\
\texttt{1: W} \\
\texttt{2: I} \\
\texttt{3: │ (boundary)} \\
\texttt{4: ◌} \\
\texttt{5: T} \\
\texttt{6: │ (boundary)} \\
\texttt{7: ◌} \\
\\ 
\texttt{So (3,1) and (6,1) are boundaries. My previous attempts were wrong.} \\
\texttt{Let's try to move T to a safe empty cell.} \\
\texttt{T is at (5,1).} \\
\texttt{(4,1) is `◌`.} \\
\\ 
\texttt{POHYB: T, (4,1)} \\
            }
        }
    }
    & & \\ \\

    \theutterance \stepcounter{utterance}  
    & & & \multicolumn{2}{p{0.3\linewidth}}{
        \cellcolor[rgb]{0.9,0.9,0.9}{
            \makecell[{{p{\linewidth}}}]{
                \texttt{\tiny{[GM$|$GM]}}
                \texttt{Invalid format: head is not empty: 'POHYB: T, (6,1)} \\
\\ 
\texttt{Wait, I see the grid again. (6,1) is ◌. Let me re{-}check the grid.} \\
\texttt{The grid at row 1 is: ║W I │ ◌ T │ ◌║ 1} \\
\texttt{Wait, the columns are:} \\
\texttt{1: W} \\
\texttt{2: I} \\
\texttt{3: │ (boundary)} \\
\texttt{4: ◌} \\
\texttt{5: T} \\
\texttt{6: │ (boundary)} \\
\texttt{7: ◌} \\
\\ 
\texttt{So (3,1) and (6,1) are boundaries. My previous attempts were wrong.} \\
\texttt{Let's try to move T to a safe empty cell.} \\
\texttt{T is at (5,1).} \\
\texttt{(4,1) is ◌.} \\
\\ 
\texttt{'} \\
            }
        }
    }
    & & \\ \\

    \theutterance \stepcounter{utterance}  
    & & \multicolumn{4}{p{0.6\linewidth}}{
        \cellcolor[rgb]{0.9,0.9,0.9}{
            \makecell[{{p{\linewidth}}}]{
                \texttt{\tiny{[P1$\langle$GM]}}
                \texttt{Trest: Vaše zpráva nesmí obsahovat nic před příkazem!} \\
\\ 
\texttt{Ujistěte se, že vaše odpověď obsahuje pouze `ŘEKNI: <ZPRÁVA>` nebo `POHYB: <OBJEKT>, (<X>, <Y>)` a nic jiného!} \\
\\ 
\texttt{Dohromady jste nasbírali 8 z 12 pokut.} \\
\texttt{Zkuste to prosím znovu!} \\
            }
        }
    }
    & \\ \\

    \theutterance \stepcounter{utterance}  
    & \multicolumn{4}{p{0.6\linewidth}}{
        \cellcolor[rgb]{0.9,0.9,0.9}{
            \makecell[{{p{\linewidth}}}]{
                \texttt{\tiny{[P1$\rangle$GM]}}
                \texttt{POHYB: T, (4,1)} \\
            }
        }
    }
    & & \\ \\

    \theutterance \stepcounter{utterance}  
    & & & \multicolumn{2}{p{0.3\linewidth}}{
        \cellcolor[rgb]{0.9,0.9,0.9}{
            \makecell[{{p{\linewidth}}}]{
                \texttt{\tiny{[GM$|$GM]}}
                \texttt{POHYB: T, (4,1)} \\
            }
        }
    }
    & & \\ \\

    \theutterance \stepcounter{utterance}  
    & & & \multicolumn{2}{p{0.3\linewidth}}{
        \cellcolor[rgb]{0.9,0.9,0.9}{
            \makecell[{{p{\linewidth}}}]{
                \texttt{\tiny{[GM$|$GM]}}
                \texttt{T moved to (4, 1)} \\
            }
        }
    }
    & & \\ \\

    \theutterance \stepcounter{utterance}  
    & & \multicolumn{4}{p{0.6\linewidth}}{
        \cellcolor[rgb]{0.9,0.9,0.9}{
            \makecell[{{p{\linewidth}}}]{
                \texttt{\tiny{[GM$\rangle$P2]}}
                \texttt{Vaše zpráva byla předána druhému hráči.} \\
\\ 
\\ 
\texttt{Právě hrajete kolo 9 z maximálně 20.} \\
\texttt{Dohromady jste nasbírali 8 z 12 pokut.} \\
\texttt{Druhý hráč přesunul objekt na své mřížce.} \\
\\ 
\\ 
\texttt{Jaký je váš další příkaz?} \\
            }
        }
    }
    & \\ \\

    \theutterance \stepcounter{utterance}  
    & & & \multicolumn{4}{p{0.6\linewidth}}{
        \cellcolor[rgb]{0.9,0.9,0.9}{
            \makecell[{{p{\linewidth}}}]{
                \texttt{\tiny{[GM$\langle$P2]}}
                \texttt{ŘEKNI: Jaký pohyb jsi udělal? Moje aktuální pozice jsou W(1,1), I(2,1), T(4,4), C(1,7), H(1,5).} \\
            }
        }
    }
     \\ \\

    \theutterance \stepcounter{utterance}  
    & & & \multicolumn{2}{p{0.3\linewidth}}{
        \cellcolor[rgb]{0.9,0.9,0.9}{
            \makecell[{{p{\linewidth}}}]{
                \texttt{\tiny{[GM$|$GM]}}
                \texttt{ŘEKNI: Jaký pohyb jsi udělal? Moje aktuální pozice jsou W(1,1), I(2,1), T(4,4), C(1,7), H(1,5).} \\
            }
        }
    }
    & & \\ \\

    \theutterance \stepcounter{utterance}  
    & & \multicolumn{4}{p{0.6\linewidth}}{
        \cellcolor[rgb]{0.9,0.9,0.9}{
            \makecell[{{p{\linewidth}}}]{
                \texttt{\tiny{[P1$\langle$GM]}}
                \texttt{Úspěšně přesunuto 'T' do (4,1). Vaše aktualizovaná mřížka vypadá takto:} \\
\texttt{```} \\
\texttt{ 1234567} \\
\texttt{╔══╤══╤═╗} \\
\texttt{║WI│T◌│◌║ 1} \\
\texttt{╟──┼──┴─╢ 2} \\
\texttt{║◌◌│◌◌◌◌║ 3} \\
\texttt{╟──┘◌◌◌◌║ 4} \\
\texttt{║◌◌◌H◌◌C║ 5} \\
\texttt{║◌◌┌─┬──╢ 6} \\
\texttt{║◌◌│◌│◌◌║ 7} \\
\texttt{╚══╧═╧══╝} \\
\texttt{```} \\
\texttt{Právě hrajete kolo 9 z maximálně 20.} \\
\texttt{Dohromady jste nasbírali 8 z 12 pokut.} \\
\texttt{Druhý hráč poslal tuto zprávu:} \\
\\ 
\\ 
\texttt{"Jaký pohyb jsi udělal? Moje aktuální pozice jsou W(1,1), I(2,1), T(4,4), C(1,7), H(1,5)."} \\
\\ 
\\ 
\texttt{Jaký je váš další příkaz?} \\
            }
        }
    }
    & \\ \\

    \theutterance \stepcounter{utterance}  
    & \multicolumn{4}{p{0.6\linewidth}}{
        \cellcolor[rgb]{0.9,0.9,0.9}{
            \makecell[{{p{\linewidth}}}]{
                \texttt{\tiny{[P1$\rangle$GM]}}
                \texttt{ŘEKNI: Přesunul jsem T na (4,1). Moje aktuální pozice jsou W(1,1), I(2,1), T(4,1), H(4,5), C(7,5). Teď se pokusím posunout T na tvé cílové místo (3,1), ale musím si dát pozor na ty čáry. Zkusím nejdřív posunout T na (4,2), pokud je to možné, nebo přímo na (4,3). Počkej, (3,1) je čára. Zkusím T posunout na (4,3).} \\
            }
        }
    }
    & & \\ \\

    \theutterance \stepcounter{utterance}  
    & & & \multicolumn{2}{p{0.3\linewidth}}{
        \cellcolor[rgb]{0.9,0.9,0.9}{
            \makecell[{{p{\linewidth}}}]{
                \texttt{\tiny{[GM$|$GM]}}
                \texttt{ŘEKNI: Přesunul jsem T na (4,1). Moje aktuální pozice jsou W(1,1), I(2,1), T(4,1), H(4,5), C(7,5). Teď se pokusím posunout T na tvé cílové místo (3,1), ale musím si dát pozor na ty čáry. Zkusím nejdřív posunout T na (4,2), pokud je to možné, nebo přímo na (4,3). Počkej, (3,1) je čára. Zkusím T posunout na (4,3).} \\
            }
        }
    }
    & & \\ \\

    \theutterance \stepcounter{utterance}  
    & & \multicolumn{4}{p{0.6\linewidth}}{
        \cellcolor[rgb]{0.9,0.9,0.9}{
            \makecell[{{p{\linewidth}}}]{
                \texttt{\tiny{[GM$\rangle$P2]}}
                \texttt{Vaše zpráva byla předána druhému hráči.} \\
\\ 
\\ 
\texttt{Právě hrajete kolo 10 z maximálně 20.} \\
\texttt{Dohromady jste nasbírali 8 z 12 pokut.} \\
\texttt{Druhý hráč poslal tuto zprávu:} \\
\\ 
\\ 
\texttt{"Přesunul jsem T na (4,1). Moje aktuální pozice jsou W(1,1), I(2,1), T(4,1), H(4,5), C(7,5). Teď se pokusím posunout T na tvé cílové místo (3,1), ale musím si dát pozor na ty čáry. Zkusím nejdřív posunout T na (4,2), pokud je to možné, nebo přímo na (4,3). Počkej, (3,1) je čára. Zkusím T posunout na (4,3)."} \\
\\ 
\\ 
\texttt{Jaký je váš další příkaz?} \\
            }
        }
    }
    & \\ \\

    \theutterance \stepcounter{utterance}  
    & & & \multicolumn{4}{p{0.6\linewidth}}{
        \cellcolor[rgb]{0.9,0.9,0.9}{
            \makecell[{{p{\linewidth}}}]{
                \texttt{\tiny{[GM$\langle$P2]}}
                \texttt{POHYB: T, (3, 1)} \\
            }
        }
    }
     \\ \\

    \theutterance \stepcounter{utterance}  
    & & & \multicolumn{2}{p{0.3\linewidth}}{
        \cellcolor[rgb]{0.9,0.9,0.9}{
            \makecell[{{p{\linewidth}}}]{
                \texttt{\tiny{[GM$|$GM]}}
                \texttt{POHYB: T, (3, 1)} \\
            }
        }
    }
    & & \\ \\

    \theutterance \stepcounter{utterance}  
    & & & \multicolumn{2}{p{0.3\linewidth}}{
        \cellcolor[rgb]{0.9,0.9,0.9}{
            \makecell[{{p{\linewidth}}}]{
                \texttt{\tiny{[GM$|$GM]}}
                \texttt{Penalizace: (3,1) není prázdný, ale obsahuje '│'.} \\
\\ 
\texttt{Dohromady jste nasbírali 9 z 12 pokut.} \\
\texttt{Zkuste to prosím znovu!} \\
            }
        }
    }
    & & \\ \\

    \theutterance \stepcounter{utterance}  
    & & \multicolumn{4}{p{0.6\linewidth}}{
        \cellcolor[rgb]{0.9,0.9,0.9}{
            \makecell[{{p{\linewidth}}}]{
                \texttt{\tiny{[GM$\rangle$P2]}}
                \texttt{Penalizace: (3,1) není prázdný, ale obsahuje '│'.} \\
\\ 
\texttt{Dohromady jste nasbírali 9 z 12 pokut.} \\
\texttt{Zkuste to prosím znovu!} \\
            }
        }
    }
    & \\ \\

    \theutterance \stepcounter{utterance}  
    & & & \multicolumn{4}{p{0.6\linewidth}}{
        \cellcolor[rgb]{0.9,0.9,0.9}{
            \makecell[{{p{\linewidth}}}]{
                \texttt{\tiny{[GM$\langle$P2]}}
                \texttt{POHYB: T, (4, 3)} \\
            }
        }
    }
     \\ \\

    \theutterance \stepcounter{utterance}  
    & & & \multicolumn{2}{p{0.3\linewidth}}{
        \cellcolor[rgb]{0.9,0.9,0.9}{
            \makecell[{{p{\linewidth}}}]{
                \texttt{\tiny{[GM$|$GM]}}
                \texttt{POHYB: T, (4, 3)} \\
            }
        }
    }
    & & \\ \\

    \theutterance \stepcounter{utterance}  
    & & & \multicolumn{2}{p{0.3\linewidth}}{
        \cellcolor[rgb]{0.9,0.9,0.9}{
            \makecell[{{p{\linewidth}}}]{
                \texttt{\tiny{[GM$|$GM]}}
                \texttt{T moved to (4, 3)} \\
            }
        }
    }
    & & \\ \\

    \theutterance \stepcounter{utterance}  
    & & \multicolumn{4}{p{0.6\linewidth}}{
        \cellcolor[rgb]{0.9,0.9,0.9}{
            \makecell[{{p{\linewidth}}}]{
                \texttt{\tiny{[P1$\langle$GM]}}
                \texttt{Vaše zpráva byla předána druhému hráči.} \\
\\ 
\\ 
\texttt{Právě hrajete kolo 10 z maximálně 20.} \\
\texttt{Dohromady jste nasbírali 9 z 12 pokut.} \\
\texttt{Druhý hráč přesunul objekt na své mřížce.} \\
\\ 
\\ 
\texttt{Jaký je váš další příkaz?} \\
            }
        }
    }
    & \\ \\

    \theutterance \stepcounter{utterance}  
    & \multicolumn{4}{p{0.6\linewidth}}{
        \cellcolor[rgb]{0.9,0.9,0.9}{
            \makecell[{{p{\linewidth}}}]{
                \texttt{\tiny{[P1$\rangle$GM]}}
                \texttt{ŘEKNI: Jaký pohyb jsi udělal? Moje aktuální pozice jsou W(1,1), I(2,1), T(4,1), H(4,5), C(7,5). Teď se pokusím posunout H na (5,1).} \\
            }
        }
    }
    & & \\ \\

    \theutterance \stepcounter{utterance}  
    & & & \multicolumn{2}{p{0.3\linewidth}}{
        \cellcolor[rgb]{0.9,0.9,0.9}{
            \makecell[{{p{\linewidth}}}]{
                \texttt{\tiny{[GM$|$GM]}}
                \texttt{ŘEKNI: Jaký pohyb jsi udělal? Moje aktuální pozice jsou W(1,1), I(2,1), T(4,1), H(4,5), C(7,5). Teď se pokusím posunout H na (5,1).} \\
            }
        }
    }
    & & \\ \\

    \theutterance \stepcounter{utterance}  
    & & \multicolumn{4}{p{0.6\linewidth}}{
        \cellcolor[rgb]{0.9,0.9,0.9}{
            \makecell[{{p{\linewidth}}}]{
                \texttt{\tiny{[GM$\rangle$P2]}}
                \texttt{Úspěšně přesunuto 'T' do (4,3). Vaše aktualizovaná mřížka vypadá takto:} \\
\texttt{```} \\
\texttt{ 1234567} \\
\texttt{╔══╤══╤═╗} \\
\texttt{║WI│◌◌│◌║ 1} \\
\texttt{╟──┼──┴─╢ 2} \\
\texttt{║◌◌│T◌◌◌║ 3} \\
\texttt{╟──┘◌◌◌◌║ 4} \\
\texttt{║H◌◌◌◌◌◌║ 5} \\
\texttt{║◌◌┌─┬──╢ 6} \\
\texttt{║C◌│◌│◌◌║ 7} \\
\texttt{╚══╧═╧══╝} \\
\texttt{```} \\
\texttt{Právě hrajete kolo 11 z maximálně 20.} \\
\texttt{Dohromady jste nasbírali 9 z 12 pokut.} \\
\texttt{Druhý hráč poslal tuto zprávu:} \\
\\ 
\\ 
\texttt{"Jaký pohyb jsi udělal? Moje aktuální pozice jsou W(1,1), I(2,1), T(4,1), H(4,5), C(7,5). Teď se pokusím posunout H na (5,1)."} \\
\\ 
\\ 
\texttt{Jaký je váš další příkaz?} \\
            }
        }
    }
    & \\ \\

    \theutterance \stepcounter{utterance}  
    & & & \multicolumn{4}{p{0.6\linewidth}}{
        \cellcolor[rgb]{0.9,0.9,0.9}{
            \makecell[{{p{\linewidth}}}]{
                \texttt{\tiny{[GM$\langle$P2]}}
                \texttt{ŘEKNI: Udělal jsem pohyb T na (4,3). Moje aktuální pozice jsou W(1,1), I(2,1), T(4,3), C(1,7), H(1,5). Teď se pokusím posunout T na (3,1), ale musím si ověřit, zda je (3,1) volné. Podívám se na mřížku... (3,1) je v mém případě prázdné, ale v tvém může být čára. Zkusím tedy nejdřív posunout H na (5,1), abychom měli volnější prostor.} \\
            }
        }
    }
     \\ \\

    \theutterance \stepcounter{utterance}  
    & & & \multicolumn{2}{p{0.3\linewidth}}{
        \cellcolor[rgb]{0.9,0.9,0.9}{
            \makecell[{{p{\linewidth}}}]{
                \texttt{\tiny{[GM$|$GM]}}
                \texttt{ŘEKNI: Udělal jsem pohyb T na (4,3). Moje aktuální pozice jsou W(1,1), I(2,1), T(4,3), C(1,7), H(1,5). Teď se pokusím posunout T na (3,1), ale musím si ověřit, zda je (3,1) volné. Podívám se na mřížku... (3,1) je v mém případě prázdné, ale v tvém může být čára. Zkusím tedy nejdřív posunout H na (5,1), abychom měli volnější prostor.} \\
            }
        }
    }
    & & \\ \\

    \theutterance \stepcounter{utterance}  
    & & \multicolumn{4}{p{0.6\linewidth}}{
        \cellcolor[rgb]{0.9,0.9,0.9}{
            \makecell[{{p{\linewidth}}}]{
                \texttt{\tiny{[P1$\langle$GM]}}
                \texttt{Vaše zpráva byla předána druhému hráči.} \\
\\ 
\\ 
\texttt{Právě hrajete kolo 11 z maximálně 20.} \\
\texttt{Dohromady jste nasbírali 9 z 12 pokut.} \\
\texttt{Druhý hráč poslal tuto zprávu:} \\
\\ 
\\ 
\texttt{"Udělal jsem pohyb T na (4,3). Moje aktuální pozice jsou W(1,1), I(2,1), T(4,3), C(1,7), H(1,5). Teď se pokusím posunout T na (3,1), ale musím si ověřit, zda je (3,1) volné. Podívám se na mřížku... (3,1) je v mém případě prázdné, ale v tvém může být čára. Zkusím tedy nejdřív posunout H na (5,1), abychom měli volnější prostor."} \\
\\ 
\\ 
\texttt{Jaký je váš další příkaz?} \\
            }
        }
    }
    & \\ \\

    \theutterance \stepcounter{utterance}  
    & \multicolumn{4}{p{0.6\linewidth}}{
        \cellcolor[rgb]{0.9,0.9,0.9}{
            \makecell[{{p{\linewidth}}}]{
                \texttt{\tiny{[P1$\rangle$GM]}}
                \texttt{POHYB: H, (5,1)} \\
            }
        }
    }
    & & \\ \\

    \theutterance \stepcounter{utterance}  
    & & & \multicolumn{2}{p{0.3\linewidth}}{
        \cellcolor[rgb]{0.9,0.9,0.9}{
            \makecell[{{p{\linewidth}}}]{
                \texttt{\tiny{[GM$|$GM]}}
                \texttt{POHYB: H, (5,1)} \\
            }
        }
    }
    & & \\ \\

    \theutterance \stepcounter{utterance}  
    & & & \multicolumn{2}{p{0.3\linewidth}}{
        \cellcolor[rgb]{0.9,0.9,0.9}{
            \makecell[{{p{\linewidth}}}]{
                \texttt{\tiny{[GM$|$GM]}}
                \texttt{H moved to (5, 1)} \\
            }
        }
    }
    & & \\ \\

    \theutterance \stepcounter{utterance}  
    & & \multicolumn{4}{p{0.6\linewidth}}{
        \cellcolor[rgb]{0.9,0.9,0.9}{
            \makecell[{{p{\linewidth}}}]{
                \texttt{\tiny{[GM$\rangle$P2]}}
                \texttt{Vaše zpráva byla předána druhému hráči.} \\
\\ 
\\ 
\texttt{Právě hrajete kolo 12 z maximálně 20.} \\
\texttt{Dohromady jste nasbírali 9 z 12 pokut.} \\
\texttt{Druhý hráč přesunul objekt na své mřížce.} \\
\\ 
\\ 
\texttt{Jaký je váš další příkaz?} \\
            }
        }
    }
    & \\ \\

    \theutterance \stepcounter{utterance}  
    & & & \multicolumn{4}{p{0.6\linewidth}}{
        \cellcolor[rgb]{0.9,0.9,0.9}{
            \makecell[{{p{\linewidth}}}]{
                \texttt{\tiny{[GM$\langle$P2]}}
                \texttt{ŘEKNI: Jaký pohyb jsi udělal? Moje aktuální pozice jsou W(1,1), I(2,1), T(4,3), C(1,7), H(1,5).} \\
            }
        }
    }
     \\ \\

    \theutterance \stepcounter{utterance}  
    & & & \multicolumn{2}{p{0.3\linewidth}}{
        \cellcolor[rgb]{0.9,0.9,0.9}{
            \makecell[{{p{\linewidth}}}]{
                \texttt{\tiny{[GM$|$GM]}}
                \texttt{ŘEKNI: Jaký pohyb jsi udělal? Moje aktuální pozice jsou W(1,1), I(2,1), T(4,3), C(1,7), H(1,5).} \\
            }
        }
    }
    & & \\ \\

    \theutterance \stepcounter{utterance}  
    & & \multicolumn{4}{p{0.6\linewidth}}{
        \cellcolor[rgb]{0.9,0.9,0.9}{
            \makecell[{{p{\linewidth}}}]{
                \texttt{\tiny{[P1$\langle$GM]}}
                \texttt{Úspěšně přesunuto 'H' do (5,1). Vaše aktualizovaná mřížka vypadá takto:} \\
\texttt{```} \\
\texttt{ 1234567} \\
\texttt{╔══╤══╤═╗} \\
\texttt{║WI│TH│◌║ 1} \\
\texttt{╟──┼──┴─╢ 2} \\
\texttt{║◌◌│◌◌◌◌║ 3} \\
\texttt{╟──┘◌◌◌◌║ 4} \\
\texttt{║◌◌◌◌◌◌C║ 5} \\
\texttt{║◌◌┌─┬──╢ 6} \\
\texttt{║◌◌│◌│◌◌║ 7} \\
\texttt{╚══╧═╧══╝} \\
\texttt{```} \\
\texttt{Právě hrajete kolo 12 z maximálně 20.} \\
\texttt{Dohromady jste nasbírali 9 z 12 pokut.} \\
\texttt{Druhý hráč poslal tuto zprávu:} \\
\\ 
\\ 
\texttt{"Jaký pohyb jsi udělal? Moje aktuální pozice jsou W(1,1), I(2,1), T(4,3), C(1,7), H(1,5)."} \\
\\ 
\\ 
\texttt{Jaký je váš další příkaz?} \\
            }
        }
    }
    & \\ \\

    \theutterance \stepcounter{utterance}  
    & \multicolumn{4}{p{0.6\linewidth}}{
        \cellcolor[rgb]{0.9,0.9,0.9}{
            \makecell[{{p{\linewidth}}}]{
                \texttt{\tiny{[P1$\rangle$GM]}}
                \texttt{ŘEKNI: Přesunul jsem H na (5,1). Moje aktuální pozice jsou W(1,1), I(2,1), T(4,1), H(5,1), C(7,5). Teď se pokusím posunout T na (3,1), ale musím si ověřit, zda je (3,1) volné. Podívám se na mřížku... (3,1) je v mém případě čára `│`. Takže T musím nejdřív posunout jinam. Zkusím T posunout na (4,2).} \\
            }
        }
    }
    & & \\ \\

    \theutterance \stepcounter{utterance}  
    & & & \multicolumn{2}{p{0.3\linewidth}}{
        \cellcolor[rgb]{0.9,0.9,0.9}{
            \makecell[{{p{\linewidth}}}]{
                \texttt{\tiny{[GM$|$GM]}}
                \texttt{ŘEKNI: Přesunul jsem H na (5,1). Moje aktuální pozice jsou W(1,1), I(2,1), T(4,1), H(5,1), C(7,5). Teď se pokusím posunout T na (3,1), ale musím si ověřit, zda je (3,1) volné. Podívám se na mřížku... (3,1) je v mém případě čára `│`. Takže T musím nejdřív posunout jinam. Zkusím T posunout na (4,2).} \\
            }
        }
    }
    & & \\ \\

    \theutterance \stepcounter{utterance}  
    & & \multicolumn{4}{p{0.6\linewidth}}{
        \cellcolor[rgb]{0.9,0.9,0.9}{
            \makecell[{{p{\linewidth}}}]{
                \texttt{\tiny{[GM$\rangle$P2]}}
                \texttt{Vaše zpráva byla předána druhému hráči.} \\
\\ 
\\ 
\texttt{Právě hrajete kolo 13 z maximálně 20.} \\
\texttt{Dohromady jste nasbírali 9 z 12 pokut.} \\
\texttt{Druhý hráč poslal tuto zprávu:} \\
\\ 
\\ 
\texttt{"Přesunul jsem H na (5,1). Moje aktuální pozice jsou W(1,1), I(2,1), T(4,1), H(5,1), C(7,5). Teď se pokusím posunout T na (3,1), ale musím si ověřit, zda je (3,1) volné. Podívám se na mřížku... (3,1) je v mém případě čára │. Takže T musím nejdřív posunout jinam. Zkusím T posunout na (4,2)."} \\
\\ 
\\ 
\texttt{Jaký je váš další příkaz?} \\
            }
        }
    }
    & \\ \\

    \theutterance \stepcounter{utterance}  
    & & & \multicolumn{4}{p{0.6\linewidth}}{
        \cellcolor[rgb]{0.9,0.9,0.9}{
            \makecell[{{p{\linewidth}}}]{
                \texttt{\tiny{[GM$\langle$P2]}}
                \texttt{POHYB: H, (5, 1)} \\
            }
        }
    }
     \\ \\

    \theutterance \stepcounter{utterance}  
    & & & \multicolumn{2}{p{0.3\linewidth}}{
        \cellcolor[rgb]{0.9,0.9,0.9}{
            \makecell[{{p{\linewidth}}}]{
                \texttt{\tiny{[GM$|$GM]}}
                \texttt{POHYB: H, (5, 1)} \\
            }
        }
    }
    & & \\ \\

    \theutterance \stepcounter{utterance}  
    & & & \multicolumn{2}{p{0.3\linewidth}}{
        \cellcolor[rgb]{0.9,0.9,0.9}{
            \makecell[{{p{\linewidth}}}]{
                \texttt{\tiny{[GM$|$GM]}}
                \texttt{H moved to (5, 1)} \\
            }
        }
    }
    & & \\ \\

    \theutterance \stepcounter{utterance}  
    & & \multicolumn{4}{p{0.6\linewidth}}{
        \cellcolor[rgb]{0.9,0.9,0.9}{
            \makecell[{{p{\linewidth}}}]{
                \texttt{\tiny{[P1$\langle$GM]}}
                \texttt{Vaše zpráva byla předána druhému hráči.} \\
\\ 
\\ 
\texttt{Právě hrajete kolo 13 z maximálně 20.} \\
\texttt{Dohromady jste nasbírali 9 z 12 pokut.} \\
\texttt{Druhý hráč přesunul objekt na své mřížce.} \\
\\ 
\\ 
\texttt{Jaký je váš další příkaz?} \\
            }
        }
    }
    & \\ \\

    \theutterance \stepcounter{utterance}  
    & \multicolumn{4}{p{0.6\linewidth}}{
        \cellcolor[rgb]{0.9,0.9,0.9}{
            \makecell[{{p{\linewidth}}}]{
                \texttt{\tiny{[P1$\rangle$GM]}}
                \texttt{ŘEKNI: Jaký pohyb jsi udělal? Moje aktuální pozice jsou W(1,1), I(2,1), T(4,1), H(5,1), C(7,5). Teď se pokusím posunout T na (4,3), aby se uvolnilo místo pro T na (3,1), pokud to bude možné.} \\
            }
        }
    }
    & & \\ \\

    \theutterance \stepcounter{utterance}  
    & & & \multicolumn{2}{p{0.3\linewidth}}{
        \cellcolor[rgb]{0.9,0.9,0.9}{
            \makecell[{{p{\linewidth}}}]{
                \texttt{\tiny{[GM$|$GM]}}
                \texttt{ŘEKNI: Jaký pohyb jsi udělal? Moje aktuální pozice jsou W(1,1), I(2,1), T(4,1), H(5,1), C(7,5). Teď se pokusím posunout T na (4,3), aby se uvolnilo místo pro T na (3,1), pokud to bude možné.} \\
            }
        }
    }
    & & \\ \\

    \theutterance \stepcounter{utterance}  
    & & \multicolumn{4}{p{0.6\linewidth}}{
        \cellcolor[rgb]{0.9,0.9,0.9}{
            \makecell[{{p{\linewidth}}}]{
                \texttt{\tiny{[GM$\rangle$P2]}}
                \texttt{Úspěšně přesunuto 'H' do (5,1). Vaše aktualizovaná mřížka vypadá takto:} \\
\texttt{```} \\
\texttt{ 1234567} \\
\texttt{╔══╤══╤═╗} \\
\texttt{║WI│◌H│◌║ 1} \\
\texttt{╟──┼──┴─╢ 2} \\
\texttt{║◌◌│T◌◌◌║ 3} \\
\texttt{╟──┘◌◌◌◌║ 4} \\
\texttt{║◌◌◌◌◌◌◌║ 5} \\
\texttt{║◌◌┌─┬──╢ 6} \\
\texttt{║C◌│◌│◌◌║ 7} \\
\texttt{╚══╧═╧══╝} \\
\texttt{```} \\
\texttt{Právě hrajete kolo 14 z maximálně 20.} \\
\texttt{Dohromady jste nasbírali 9 z 12 pokut.} \\
\texttt{Druhý hráč poslal tuto zprávu:} \\
\\ 
\\ 
\texttt{"Jaký pohyb jsi udělal? Moje aktuální pozice jsou W(1,1), I(2,1), T(4,1), H(5,1), C(7,5). Teď se pokusím posunout T na (4,3), aby se uvolnilo místo pro T na (3,1), pokud to bude možné."} \\
\\ 
\\ 
\texttt{Jaký je váš další příkaz?} \\
            }
        }
    }
    & \\ \\

    \theutterance \stepcounter{utterance}  
    & & & \multicolumn{4}{p{0.6\linewidth}}{
        \cellcolor[rgb]{0.9,0.9,0.9}{
            \makecell[{{p{\linewidth}}}]{
                \texttt{\tiny{[GM$\langle$P2]}}
                \texttt{ŘEKNI: Udělal jsem pohyb H na (5,1). Moje aktuální pozice jsou W(1,1), I(2,1), T(4,3), H(5,1), C(1,7). Teď se pokusím posunout T na (3,1), pokud to bude možné.} \\
            }
        }
    }
     \\ \\

    \theutterance \stepcounter{utterance}  
    & & & \multicolumn{2}{p{0.3\linewidth}}{
        \cellcolor[rgb]{0.9,0.9,0.9}{
            \makecell[{{p{\linewidth}}}]{
                \texttt{\tiny{[GM$|$GM]}}
                \texttt{ŘEKNI: Udělal jsem pohyb H na (5,1). Moje aktuální pozice jsou W(1,1), I(2,1), T(4,3), H(5,1), C(1,7). Teď se pokusím posunout T na (3,1), pokud to bude možné.} \\
            }
        }
    }
    & & \\ \\

    \theutterance \stepcounter{utterance}  
    & & \multicolumn{4}{p{0.6\linewidth}}{
        \cellcolor[rgb]{0.9,0.9,0.9}{
            \makecell[{{p{\linewidth}}}]{
                \texttt{\tiny{[P1$\langle$GM]}}
                \texttt{Vaše zpráva byla předána druhému hráči.} \\
\\ 
\\ 
\texttt{Právě hrajete kolo 14 z maximálně 20.} \\
\texttt{Dohromady jste nasbírali 9 z 12 pokut.} \\
\texttt{Druhý hráč poslal tuto zprávu:} \\
\\ 
\\ 
\texttt{"Udělal jsem pohyb H na (5,1). Moje aktuální pozice jsou W(1,1), I(2,1), T(4,3), H(5,1), C(1,7). Teď se pokusím posunout T na (3,1), pokud to bude možné."} \\
\\ 
\\ 
\texttt{Jaký je váš další příkaz?} \\
            }
        }
    }
    & \\ \\

    \theutterance \stepcounter{utterance}  
    & \multicolumn{4}{p{0.6\linewidth}}{
        \cellcolor[rgb]{0.9,0.9,0.9}{
            \makecell[{{p{\linewidth}}}]{
                \texttt{\tiny{[P1$\rangle$GM]}}
                \texttt{POHYB: T, (4,3)} \\
            }
        }
    }
    & & \\ \\

    \theutterance \stepcounter{utterance}  
    & & & \multicolumn{2}{p{0.3\linewidth}}{
        \cellcolor[rgb]{0.9,0.9,0.9}{
            \makecell[{{p{\linewidth}}}]{
                \texttt{\tiny{[GM$|$GM]}}
                \texttt{POHYB: T, (4,3)} \\
            }
        }
    }
    & & \\ \\

    \theutterance \stepcounter{utterance}  
    & & & \multicolumn{2}{p{0.3\linewidth}}{
        \cellcolor[rgb]{0.9,0.9,0.9}{
            \makecell[{{p{\linewidth}}}]{
                \texttt{\tiny{[GM$|$GM]}}
                \texttt{T moved to (4, 3)} \\
            }
        }
    }
    & & \\ \\

    \theutterance \stepcounter{utterance}  
    & & \multicolumn{4}{p{0.6\linewidth}}{
        \cellcolor[rgb]{0.9,0.9,0.9}{
            \makecell[{{p{\linewidth}}}]{
                \texttt{\tiny{[GM$\rangle$P2]}}
                \texttt{Vaše zpráva byla předána druhému hráči.} \\
\\ 
\\ 
\texttt{Právě hrajete kolo 15 z maximálně 20.} \\
\texttt{Dohromady jste nasbírali 9 z 12 pokut.} \\
\texttt{Druhý hráč přesunul objekt na své mřížce.} \\
\\ 
\\ 
\texttt{Jaký je váš další příkaz?} \\
            }
        }
    }
    & \\ \\

    \theutterance \stepcounter{utterance}  
    & & & \multicolumn{4}{p{0.6\linewidth}}{
        \cellcolor[rgb]{0.9,0.9,0.9}{
            \makecell[{{p{\linewidth}}}]{
                \texttt{\tiny{[GM$\langle$P2]}}
                \texttt{ŘEKNI: Jaký pohyb jsi udělal? Moje aktuální pozice jsou W(1,1), I(2,1), T(4,3), H(5,1), C(1,7).} \\
            }
        }
    }
     \\ \\

    \theutterance \stepcounter{utterance}  
    & & & \multicolumn{2}{p{0.3\linewidth}}{
        \cellcolor[rgb]{0.9,0.9,0.9}{
            \makecell[{{p{\linewidth}}}]{
                \texttt{\tiny{[GM$|$GM]}}
                \texttt{ŘEKNI: Jaký pohyb jsi udělal? Moje aktuální pozice jsou W(1,1), I(2,1), T(4,3), H(5,1), C(1,7).} \\
            }
        }
    }
    & & \\ \\

    \theutterance \stepcounter{utterance}  
    & & \multicolumn{4}{p{0.6\linewidth}}{
        \cellcolor[rgb]{0.9,0.9,0.9}{
            \makecell[{{p{\linewidth}}}]{
                \texttt{\tiny{[P1$\langle$GM]}}
                \texttt{Úspěšně přesunuto 'T' do (4,3). Vaše aktualizovaná mřížka vypadá takto:} \\
\texttt{```} \\
\texttt{ 1234567} \\
\texttt{╔══╤══╤═╗} \\
\texttt{║WI│◌H│◌║ 1} \\
\texttt{╟──┼──┴─╢ 2} \\
\texttt{║◌◌│T◌◌◌║ 3} \\
\texttt{╟──┘◌◌◌◌║ 4} \\
\texttt{║◌◌◌◌◌◌C║ 5} \\
\texttt{║◌◌┌─┬──╢ 6} \\
\texttt{║◌◌│◌│◌◌║ 7} \\
\texttt{╚══╧═╧══╝} \\
\texttt{```} \\
\texttt{Právě hrajete kolo 15 z maximálně 20.} \\
\texttt{Dohromady jste nasbírali 9 z 12 pokut.} \\
\texttt{Druhý hráč poslal tuto zprávu:} \\
\\ 
\\ 
\texttt{"Jaký pohyb jsi udělal? Moje aktuální pozice jsou W(1,1), I(2,1), T(4,3), H(5,1), C(1,7)."} \\
\\ 
\\ 
\texttt{Jaký je váš další příkaz?} \\
            }
        }
    }
    & \\ \\

    \theutterance \stepcounter{utterance}  
    & \multicolumn{4}{p{0.6\linewidth}}{
        \cellcolor[rgb]{0.9,0.9,0.9}{
            \makecell[{{p{\linewidth}}}]{
                \texttt{\tiny{[P1$\rangle$GM]}}
                \texttt{ŘEKNI: Přesunul jsem T na (4,3). Moje aktuální pozice jsou W(1,1), I(2,1), T(4,3), H(5,1), C(7,5). Teď se pokusím posunout C na (4,1), aby se to přiblížilo tvému cíli.} \\
            }
        }
    }
    & & \\ \\

    \theutterance \stepcounter{utterance}  
    & & & \multicolumn{2}{p{0.3\linewidth}}{
        \cellcolor[rgb]{0.9,0.9,0.9}{
            \makecell[{{p{\linewidth}}}]{
                \texttt{\tiny{[GM$|$GM]}}
                \texttt{ŘEKNI: Přesunul jsem T na (4,3). Moje aktuální pozice jsou W(1,1), I(2,1), T(4,3), H(5,1), C(7,5). Teď se pokusím posunout C na (4,1), aby se to přiblížilo tvému cíli.} \\
            }
        }
    }
    & & \\ \\

    \theutterance \stepcounter{utterance}  
    & & \multicolumn{4}{p{0.6\linewidth}}{
        \cellcolor[rgb]{0.9,0.9,0.9}{
            \makecell[{{p{\linewidth}}}]{
                \texttt{\tiny{[GM$\rangle$P2]}}
                \texttt{Vaše zpráva byla předána druhému hráči.} \\
\\ 
\\ 
\texttt{Právě hrajete kolo 16 z maximálně 20.} \\
\texttt{Dohromady jste nasbírali 9 z 12 pokut.} \\
\texttt{Druhý hráč poslal tuto zprávu:} \\
\\ 
\\ 
\texttt{"Přesunul jsem T na (4,3). Moje aktuální pozice jsou W(1,1), I(2,1), T(4,3), H(5,1), C(7,5). Teď se pokusím posunout C na (4,1), aby se to přiblížilo tvému cíli."} \\
\\ 
\\ 
\texttt{Jaký je váš další příkaz?} \\
            }
        }
    }
    & \\ \\

    \theutterance \stepcounter{utterance}  
    & & & \multicolumn{4}{p{0.6\linewidth}}{
        \cellcolor[rgb]{0.9,0.9,0.9}{
            \makecell[{{p{\linewidth}}}]{
                \texttt{\tiny{[GM$\langle$P2]}}
                \texttt{POHYB: C, (4, 1)} \\
            }
        }
    }
     \\ \\

    \theutterance \stepcounter{utterance}  
    & & & \multicolumn{2}{p{0.3\linewidth}}{
        \cellcolor[rgb]{0.9,0.9,0.9}{
            \makecell[{{p{\linewidth}}}]{
                \texttt{\tiny{[GM$|$GM]}}
                \texttt{POHYB: C, (4, 1)} \\
            }
        }
    }
    & & \\ \\

    \theutterance \stepcounter{utterance}  
    & & & \multicolumn{2}{p{0.3\linewidth}}{
        \cellcolor[rgb]{0.9,0.9,0.9}{
            \makecell[{{p{\linewidth}}}]{
                \texttt{\tiny{[GM$|$GM]}}
                \texttt{C moved to (4, 1)} \\
            }
        }
    }
    & & \\ \\

    \theutterance \stepcounter{utterance}  
    & & \multicolumn{4}{p{0.6\linewidth}}{
        \cellcolor[rgb]{0.9,0.9,0.9}{
            \makecell[{{p{\linewidth}}}]{
                \texttt{\tiny{[P1$\langle$GM]}}
                \texttt{Vaše zpráva byla předána druhému hráči.} \\
\\ 
\\ 
\texttt{Právě hrajete kolo 16 z maximálně 20.} \\
\texttt{Dohromady jste nasbírali 9 z 12 pokut.} \\
\texttt{Druhý hráč přesunul objekt na své mřížce.} \\
\\ 
\\ 
\texttt{Jaký je váš další příkaz?} \\
            }
        }
    }
    & \\ \\

    \theutterance \stepcounter{utterance}  
    & \multicolumn{4}{p{0.6\linewidth}}{
        \cellcolor[rgb]{0.9,0.9,0.9}{
            \makecell[{{p{\linewidth}}}]{
                \texttt{\tiny{[P1$\rangle$GM]}}
                \texttt{ŘEKNI: Jaký pohyb jsi udělal? Moje aktuální pozice jsou W(1,1), I(2,1), T(4,3), H(5,1), C(7,5). Teď se pokusím posunout C na (4,1).} \\
            }
        }
    }
    & & \\ \\

    \theutterance \stepcounter{utterance}  
    & & & \multicolumn{2}{p{0.3\linewidth}}{
        \cellcolor[rgb]{0.9,0.9,0.9}{
            \makecell[{{p{\linewidth}}}]{
                \texttt{\tiny{[GM$|$GM]}}
                \texttt{ŘEKNI: Jaký pohyb jsi udělal? Moje aktuální pozice jsou W(1,1), I(2,1), T(4,3), H(5,1), C(7,5). Teď se pokusím posunout C na (4,1).} \\
            }
        }
    }
    & & \\ \\

    \theutterance \stepcounter{utterance}  
    & & \multicolumn{4}{p{0.6\linewidth}}{
        \cellcolor[rgb]{0.9,0.9,0.9}{
            \makecell[{{p{\linewidth}}}]{
                \texttt{\tiny{[GM$\rangle$P2]}}
                \texttt{Úspěšně přesunuto 'C' do (4,1). Vaše aktualizovaná mřížka vypadá takto:} \\
\texttt{```} \\
\texttt{ 1234567} \\
\texttt{╔══╤══╤═╗} \\
\texttt{║WI│CH│◌║ 1} \\
\texttt{╟──┼──┴─╢ 2} \\
\texttt{║◌◌│T◌◌◌║ 3} \\
\texttt{╟──┘◌◌◌◌║ 4} \\
\texttt{║◌◌◌◌◌◌◌║ 5} \\
\texttt{║◌◌┌─┬──╢ 6} \\
\texttt{║◌◌│◌│◌◌║ 7} \\
\texttt{╚══╧═╧══╝} \\
\texttt{```} \\
\texttt{Právě hrajete kolo 17 z maximálně 20.} \\
\texttt{Dohromady jste nasbírali 9 z 12 pokut.} \\
\texttt{Druhý hráč poslal tuto zprávu:} \\
\\ 
\\ 
\texttt{"Jaký pohyb jsi udělal? Moje aktuální pozice jsou W(1,1), I(2,1), T(4,3), H(5,1), C(7,5). Teď se pokusím posunout C na (4,1)."} \\
\\ 
\\ 
\texttt{Jaký je váš další příkaz?} \\
            }
        }
    }
    & \\ \\

    \theutterance \stepcounter{utterance}  
    & & & \multicolumn{4}{p{0.6\linewidth}}{
        \cellcolor[rgb]{0.9,0.9,0.9}{
            \makecell[{{p{\linewidth}}}]{
                \texttt{\tiny{[GM$\langle$P2]}}
                \texttt{ŘEKNI: Udělal jsem pohyb C na (4,1). Moje aktuální pozice jsou W(1,1), I(2,1), T(4,3), H(5,1), C(4,1). Teď už zbývá jen T. Moje T je na (4,3), tvé T je na (4,3). Musíme se shodnout, kde má T být. Navrhuji (3,1), ale musíme si ověřit, zda tam není překážka. Pokud ano, navrhuji jinou souřadnici. Jakou máš teď C?} \\
            }
        }
    }
     \\ \\

    \theutterance \stepcounter{utterance}  
    & & & \multicolumn{2}{p{0.3\linewidth}}{
        \cellcolor[rgb]{0.9,0.9,0.9}{
            \makecell[{{p{\linewidth}}}]{
                \texttt{\tiny{[GM$|$GM]}}
                \texttt{ŘEKNI: Udělal jsem pohyb C na (4,1). Moje aktuální pozice jsou W(1,1), I(2,1), T(4,3), H(5,1), C(4,1). Teď už zbývá jen T. Moje T je na (4,3), tvé T je na (4,3). Musíme se shodnout, kde má T být. Navrhuji (3,1), ale musíme si ověřit, zda tam není překážka. Pokud ano, navrhuji jinou souřadnici. Jakou máš teď C?} \\
            }
        }
    }
    & & \\ \\

    \theutterance \stepcounter{utterance}  
    & & \multicolumn{4}{p{0.6\linewidth}}{
        \cellcolor[rgb]{0.9,0.9,0.9}{
            \makecell[{{p{\linewidth}}}]{
                \texttt{\tiny{[P1$\langle$GM]}}
                \texttt{Vaše zpráva byla předána druhému hráči.} \\
\\ 
\\ 
\texttt{Právě hrajete kolo 17 z maximálně 20.} \\
\texttt{Dohromady jste nasbírali 9 z 12 pokut.} \\
\texttt{Druhý hráč poslal tuto zprávu:} \\
\\ 
\\ 
\texttt{"Udělal jsem pohyb C na (4,1). Moje aktuální pozice jsou W(1,1), I(2,1), T(4,3), H(5,1), C(4,1). Teď už zbývá jen T. Moje T je na (4,3), tvé T je na (4,3). Musíme se shodnout, kde má T být. Navrhuji (3,1), ale musíme si ověřit, zda tam není překážka. Pokud ano, navrhuji jinou souřadnici. Jakou máš teď C?"} \\
\\ 
\\ 
\texttt{Jaký je váš další příkaz?} \\
            }
        }
    }
    & \\ \\

    \theutterance \stepcounter{utterance}  
    & \multicolumn{4}{p{0.6\linewidth}}{
        \cellcolor[rgb]{0.9,0.9,0.9}{
            \makecell[{{p{\linewidth}}}]{
                \texttt{\tiny{[P1$\rangle$GM]}}
                \texttt{POHYB: C, (4,1)} \\
            }
        }
    }
    & & \\ \\

    \theutterance \stepcounter{utterance}  
    & & & \multicolumn{2}{p{0.3\linewidth}}{
        \cellcolor[rgb]{0.9,0.9,0.9}{
            \makecell[{{p{\linewidth}}}]{
                \texttt{\tiny{[GM$|$GM]}}
                \texttt{POHYB: C, (4,1)} \\
            }
        }
    }
    & & \\ \\

    \theutterance \stepcounter{utterance}  
    & & & \multicolumn{2}{p{0.3\linewidth}}{
        \cellcolor[rgb]{0.9,0.9,0.9}{
            \makecell[{{p{\linewidth}}}]{
                \texttt{\tiny{[GM$|$GM]}}
                \texttt{C moved to (4, 1)} \\
            }
        }
    }
    & & \\ \\

    \theutterance \stepcounter{utterance}  
    & & \multicolumn{4}{p{0.6\linewidth}}{
        \cellcolor[rgb]{0.9,0.9,0.9}{
            \makecell[{{p{\linewidth}}}]{
                \texttt{\tiny{[GM$\rangle$P2]}}
                \texttt{Vaše zpráva byla předána druhému hráči.} \\
\\ 
\\ 
\texttt{Právě hrajete kolo 18 z maximálně 20.} \\
\texttt{Dohromady jste nasbírali 9 z 12 pokut.} \\
\texttt{Druhý hráč přesunul objekt na své mřížce.} \\
\\ 
\\ 
\texttt{Jaký je váš další příkaz?} \\
            }
        }
    }
    & \\ \\

    \theutterance \stepcounter{utterance}  
    & & & \multicolumn{4}{p{0.6\linewidth}}{
        \cellcolor[rgb]{0.9,0.9,0.9}{
            \makecell[{{p{\linewidth}}}]{
                \texttt{\tiny{[GM$\langle$P2]}}
                \texttt{ŘEKNI: Jaký pohyb jsi udělal? Moje aktuální pozice jsou W(1,1), I(2,1), T(4,3), H(5,1), C(4,1). Pokud jsi posunul C na (4,1), tak jsme téměř hotovi, jen zbývá T.} \\
            }
        }
    }
     \\ \\

    \theutterance \stepcounter{utterance}  
    & & & \multicolumn{2}{p{0.3\linewidth}}{
        \cellcolor[rgb]{0.9,0.9,0.9}{
            \makecell[{{p{\linewidth}}}]{
                \texttt{\tiny{[GM$|$GM]}}
                \texttt{ŘEKNI: Jaký pohyb jsi udělal? Moje aktuální pozice jsou W(1,1), I(2,1), T(4,3), H(5,1), C(4,1). Pokud jsi posunul C na (4,1), tak jsme téměř hotovi, jen zbývá T.} \\
            }
        }
    }
    & & \\ \\

    \theutterance \stepcounter{utterance}  
    & & \multicolumn{4}{p{0.6\linewidth}}{
        \cellcolor[rgb]{0.9,0.9,0.9}{
            \makecell[{{p{\linewidth}}}]{
                \texttt{\tiny{[P1$\langle$GM]}}
                \texttt{Úspěšně přesunuto 'C' do (4,1). Vaše aktualizovaná mřížka vypadá takto:} \\
\texttt{```} \\
\texttt{ 1234567} \\
\texttt{╔══╤══╤═╗} \\
\texttt{║WI│CH│◌║ 1} \\
\texttt{╟──┼──┴─╢ 2} \\
\texttt{║◌◌│T◌◌◌║ 3} \\
\texttt{╟──┘◌◌◌◌║ 4} \\
\texttt{║◌◌◌◌◌◌◌║ 5} \\
\texttt{║◌◌┌─┬──╢ 6} \\
\texttt{║◌◌│◌│◌◌║ 7} \\
\texttt{╚══╧═╧══╝} \\
\texttt{```} \\
\texttt{Právě hrajete kolo 18 z maximálně 20.} \\
\texttt{Dohromady jste nasbírali 9 z 12 pokut.} \\
\texttt{Druhý hráč poslal tuto zprávu:} \\
\\ 
\\ 
\texttt{"Jaký pohyb jsi udělal? Moje aktuální pozice jsou W(1,1), I(2,1), T(4,3), H(5,1), C(4,1). Pokud jsi posunul C na (4,1), tak jsme téměř hotovi, jen zbývá T."} \\
\\ 
\\ 
\texttt{Jaký je váš další příkaz?} \\
            }
        }
    }
    & \\ \\

    \theutterance \stepcounter{utterance}  
    & \multicolumn{4}{p{0.6\linewidth}}{
        \cellcolor[rgb]{0.9,0.9,0.9}{
            \makecell[{{p{\linewidth}}}]{
                \texttt{\tiny{[P1$\rangle$GM]}}
                \texttt{ŘEKNI: Ano, posunul jsem C na (4,1). Moje aktuální pozice jsou W(1,1), I(2,1), T(4,3), H(5,1), C(4,1). Tvá pozice T je (4,3) a moje také (4,3). To je skvělé! Teď musíme vyřešit T. Protože (3,1) je u mě čára, navrhuji, aby T zůstalo na (4,3) nebo se přesunulo na (4,2), pokud je to u tebe volné. Ale počkej, tvůj původní plán byl T(3,1). Pokud (3,1) u mě není volné, navrhuji, aby T skončilo na (4,3) pro oba. Co ty na to? Pokud souhlasíš, můžeme zkusit skončit.} \\
            }
        }
    }
    & & \\ \\

    \theutterance \stepcounter{utterance}  
    & & & \multicolumn{2}{p{0.3\linewidth}}{
        \cellcolor[rgb]{0.9,0.9,0.9}{
            \makecell[{{p{\linewidth}}}]{
                \texttt{\tiny{[GM$|$GM]}}
                \texttt{ŘEKNI: Ano, posunul jsem C na (4,1). Moje aktuální pozice jsou W(1,1), I(2,1), T(4,3), H(5,1), C(4,1). Tvá pozice T je (4,3) a moje také (4,3). To je skvělé! Teď musíme vyřešit T. Protože (3,1) je u mě čára, navrhuji, aby T zůstalo na (4,3) nebo se přesunulo na (4,2), pokud je to u tebe volné. Ale počkej, tvůj původní plán byl T(3,1). Pokud (3,1) u mě není volné, navrhuji, aby T skončilo na (4,3) pro oba. Co ty na to? Pokud souhlasíš, můžeme zkusit skončit.} \\
            }
        }
    }
    & & \\ \\

    \theutterance \stepcounter{utterance}  
    & & \multicolumn{4}{p{0.6\linewidth}}{
        \cellcolor[rgb]{0.9,0.9,0.9}{
            \makecell[{{p{\linewidth}}}]{
                \texttt{\tiny{[GM$\rangle$P2]}}
                \texttt{Vaše zpráva byla předána druhému hráči.} \\
\\ 
\\ 
\texttt{Právě hrajete kolo 19 z maximálně 20.} \\
\texttt{Dohromady jste nasbírali 9 z 12 pokut.} \\
\texttt{Druhý hráč poslal tuto zprávu:} \\
\\ 
\\ 
\texttt{"Ano, posunul jsem C na (4,1). Moje aktuální pozice jsou W(1,1), I(2,1), T(4,3), H(5,1), C(4,1). Tvá pozice T je (4,3) a moje také (4,3). To je skvělé! Teď musíme vyřešit T. Protože (3,1) je u mě čára, navrhuji, aby T zůstalo na (4,3) nebo se přesunulo na (4,2), pokud je to u tebe volné. Ale počkej, tvůj původní plán byl T(3,1). Pokud (3,1) u mě není volné, navrhuji, aby T skončilo na (4,3) pro oba. Co ty na to? Pokud souhlasíš, můžeme zkusit skončit."} \\
\\ 
\\ 
\texttt{Jaký je váš další příkaz?} \\
            }
        }
    }
    & \\ \\

    \theutterance \stepcounter{utterance}  
    & & & \multicolumn{4}{p{0.6\linewidth}}{
        \cellcolor[rgb]{0.9,0.9,0.9}{
            \makecell[{{p{\linewidth}}}]{
                \texttt{\tiny{[GM$\langle$P2]}}
                \texttt{ŘEKNI: Souhlasím, T(4,3) je skvělé řešení, protože to u nás obou funguje a je to bezpečné. Všechny objekty jsou na stejných pozicích: W(1,1), I(2,1), T(4,3), C(4,1), H(5,1). Hotovo?} \\
            }
        }
    }
     \\ \\

    \theutterance \stepcounter{utterance}  
    & & & \multicolumn{2}{p{0.3\linewidth}}{
        \cellcolor[rgb]{0.9,0.9,0.9}{
            \makecell[{{p{\linewidth}}}]{
                \texttt{\tiny{[GM$|$GM]}}
                \texttt{ŘEKNI: Souhlasím, T(4,3) je skvělé řešení, protože to u nás obou funguje a je to bezpečné. Všechny objekty jsou na stejných pozicích: W(1,1), I(2,1), T(4,3), C(4,1), H(5,1). Hotovo?} \\
            }
        }
    }
    & & \\ \\

    \theutterance \stepcounter{utterance}  
    & & \multicolumn{4}{p{0.6\linewidth}}{
        \cellcolor[rgb]{0.9,0.9,0.9}{
            \makecell[{{p{\linewidth}}}]{
                \texttt{\tiny{[P1$\langle$GM]}}
                \texttt{Vaše zpráva byla předána druhému hráči.} \\
\\ 
\\ 
\texttt{Právě hrajete kolo 19 z maximálně 20.} \\
\texttt{Dohromady jste nasbírali 9 z 12 pokut.} \\
\texttt{Druhý hráč poslal tuto zprávu:} \\
\\ 
\\ 
\texttt{"Souhlasím, T(4,3) je skvělé řešení, protože to u nás obou funguje a je to bezpečné. Všechny objekty jsou na stejných pozicích: W(1,1), I(2,1), T(4,3), C(4,1), H(5,1). Hotovo?"} \\
\\ 
\\ 
\texttt{Jaký je váš další příkaz?} \\
            }
        }
    }
    & \\ \\

    \theutterance \stepcounter{utterance}  
    & \multicolumn{4}{p{0.6\linewidth}}{
        \cellcolor[rgb]{0.9,0.9,0.9}{
            \makecell[{{p{\linewidth}}}]{
                \texttt{\tiny{[P1$\rangle$GM]}}
                \texttt{ŘEKNI: hotovo!} \\
            }
        }
    }
    & & \\ \\

    \theutterance \stepcounter{utterance}  
    & & & \multicolumn{2}{p{0.3\linewidth}}{
        \cellcolor[rgb]{0.9,0.9,0.9}{
            \makecell[{{p{\linewidth}}}]{
                \texttt{\tiny{[GM$|$GM]}}
                \texttt{ŘEKNI: hotovo!} \\
            }
        }
    }
    & & \\ \\

    \theutterance \stepcounter{utterance}  
    & & \multicolumn{4}{p{0.6\linewidth}}{
        \cellcolor[rgb]{0.9,0.9,0.9}{
            \makecell[{{p{\linewidth}}}]{
                \texttt{\tiny{[GM$\rangle$P2]}}
                \texttt{Vaše zpráva byla předána druhému hráči.} \\
\\ 
\\ 
\texttt{Právě hrajete kolo 20 z maximálně 20.} \\
\texttt{Dohromady jste nasbírali 9 z 12 pokut.} \\
\texttt{Druhý hráč poslal tuto zprávu:} \\
\\ 
\\ 
\texttt{"hotovo!"} \\
\\ 
\\ 
\texttt{Jaký je váš další příkaz?} \\
            }
        }
    }
    & \\ \\

    \theutterance \stepcounter{utterance}  
    & & & \multicolumn{4}{p{0.6\linewidth}}{
        \cellcolor[rgb]{0.9,0.9,0.9}{
            \makecell[{{p{\linewidth}}}]{
                \texttt{\tiny{[GM$\langle$P2]}}
                \texttt{ŘEKNI: hotovo!} \\
            }
        }
    }
     \\ \\

    \theutterance \stepcounter{utterance}  
    & & & \multicolumn{2}{p{0.3\linewidth}}{
        \cellcolor[rgb]{0.9,0.9,0.9}{
            \makecell[{{p{\linewidth}}}]{
                \texttt{\tiny{[GM$|$GM]}}
                \texttt{ŘEKNI: hotovo!} \\
            }
        }
    }
    & & \\ \\

    \theutterance \stepcounter{utterance}  
    & & & \multicolumn{2}{p{0.3\linewidth}}{
        \cellcolor[rgb]{0.9,0.9,0.9}{
            \makecell[{{p{\linewidth}}}]{
                \texttt{\tiny{[GM$|$GM]}}
                \texttt{true} \\
            }
        }
    }
    & & \\ \\

    \theutterance \stepcounter{utterance}  
    & & & \multicolumn{2}{p{0.3\linewidth}}{
        \cellcolor[rgb]{0.9,0.9,0.9}{
            \makecell[{{p{\linewidth}}}]{
                \texttt{\tiny{[GM$|$GM]}}
                \texttt{Maximum number of rounds reached} \\
            }
        }
    }
    & & \\ \\

    \theutterance \stepcounter{utterance}  
    & & & \multicolumn{2}{p{0.3\linewidth}}{
        \cellcolor[rgb]{0.9,0.9,0.9}{
            \makecell[{{p{\linewidth}}}]{
                \texttt{\tiny{[GM$|$GM]}}
                \texttt{Initial states:} \\
\texttt{```} \\
\texttt{Player 1:} \\
\texttt{ 1234567} \\
\texttt{╔══╤══╤═╗} \\
\texttt{║◌◌│T◌│◌║ 1} \\
\texttt{╟──┼──┴─╢ 2} \\
\texttt{║WI│◌◌◌◌║ 3} \\
\texttt{╟──┘◌◌◌◌║ 4} \\
\texttt{║◌◌◌H◌◌C║ 5} \\
\texttt{║◌◌┌─┬──╢ 6} \\
\texttt{║◌◌│◌│◌◌║ 7} \\
\texttt{╚══╧═╧══╝} \\
\\ 
\\ 
\texttt{Player 2:} \\
\texttt{ 1234567} \\
\texttt{╔══╤══╤═╗} \\
\texttt{║◌◌│W◌│I║ 1} \\
\texttt{╟──┼──┴─╢ 2} \\
\texttt{║◌◌│◌◌◌◌║ 3} \\
\texttt{╟──┘T◌◌◌║ 4} \\
\texttt{║H◌◌◌◌◌◌║ 5} \\
\texttt{║◌◌┌─┬──╢ 6} \\
\texttt{║C◌│◌│◌◌║ 7} \\
\texttt{╚══╧═╧══╝} \\
\\ 
\texttt{```} \\
            }
        }
    }
    & & \\ \\

    \theutterance \stepcounter{utterance}  
    & & & \multicolumn{2}{p{0.3\linewidth}}{
        \cellcolor[rgb]{0.9,0.9,0.9}{
            \makecell[{{p{\linewidth}}}]{
                \texttt{\tiny{[GM$|$GM]}}
                \texttt{End states:} \\
\texttt{```} \\
\texttt{Player 1:} \\
\texttt{ 1234567} \\
\texttt{╔══╤══╤═╗} \\
\texttt{║WI│CH│◌║ 1} \\
\texttt{╟──┼──┴─╢ 2} \\
\texttt{║◌◌│T◌◌◌║ 3} \\
\texttt{╟──┘◌◌◌◌║ 4} \\
\texttt{║◌◌◌◌◌◌◌║ 5} \\
\texttt{║◌◌┌─┬──╢ 6} \\
\texttt{║◌◌│◌│◌◌║ 7} \\
\texttt{╚══╧═╧══╝} \\
\\ 
\\ 
\texttt{Player 2:} \\
\texttt{ 1234567} \\
\texttt{╔══╤══╤═╗} \\
\texttt{║WI│CH│◌║ 1} \\
\texttt{╟──┼──┴─╢ 2} \\
\texttt{║◌◌│T◌◌◌║ 3} \\
\texttt{╟──┘◌◌◌◌║ 4} \\
\texttt{║◌◌◌◌◌◌◌║ 5} \\
\texttt{║◌◌┌─┬──╢ 6} \\
\texttt{║◌◌│◌│◌◌║ 7} \\
\texttt{╚══╧═╧══╝} \\
\\ 
\texttt{```} \\
            }
        }
    }
    & & \\ \\

    \theutterance \stepcounter{utterance}  
    & & & \multicolumn{2}{p{0.3\linewidth}}{
        \cellcolor[rgb]{0.9,0.9,0.9}{
            \makecell[{{p{\linewidth}}}]{
                \texttt{\tiny{[GM$|$GM]}}
                \texttt{* success: True} \\
\texttt{* lose: False} \\
\texttt{* aborted: False} \\
\texttt{{-}{-}{-}{-}{-}{-}{-}} \\
\texttt{* Shifts: 6.00} \\
\texttt{* Max Shifts: 8.00} \\
\texttt{* Min Shifts: 4.00} \\
\texttt{* End Distance Sum: 0.00} \\
\texttt{* Init Distance Sum: 21.32} \\
\texttt{* Expected Distance Sum: 20.95} \\
\texttt{* Penalties: 9.00} \\
\texttt{* Max Penalties: 12.00} \\
\texttt{* Rounds: 20.00} \\
\texttt{* Max Rounds: 20.00} \\
\texttt{* Object Count: 5.00} \\
            }
        }
    }
    & & \\ \\

\end{supertabular}
}

\end{document}
